% concatenated from paper_arxive.tex
%\documentclass[letterpaper, 10 pt, conference]{ieeeconf}  % Comment this line out if you need a4paper
\documentclass[journal,twoside,web]{ieeecolor}

%\IEEEoverridecommandlockouts                              % This command is only needed if
                                                          % you want to use the \thanks command

%\overrideIEEEmargins      

\usepackage{amsmath}
%\usepackage{amsthm}
\usepackage{amssymb}
\usepackage{mathtools}
\usepackage{enumerate}
\usepackage{graphicx}
\usepackage{textcomp}
%\usepackage{graphicx}      % include this line if your document contains figures
%\usepackage{natbib}          % required for bibliography

% The following packages can be found on http:\\www.ctan.org
\usepackage{algorithm}
\usepackage{algorithmicx}
\usepackage{cite}
\usepackage{mathrsfs}
\usepackage{lcsys}


%\theoremstyle{plain}
\newtheorem{Theorem}{Theorem}[section]
\newtheorem{Notation}{Notation}[section]
\newtheorem{Assumption}{Assumption}[section]
\newtheorem{Proposition}[Theorem]{Proposition}
\newtheorem{Lemma}[Theorem]{Lemma}
\newtheorem{Corollary}[Theorem]{Corollary}
\newtheorem{corollary}[Theorem]{Corollary}
%\theoremstyle{definition}
\newtheorem{Definition}[Theorem]{Definition}
\newtheorem{assumption}[Theorem]{Assumption}
\newtheorem{notation}{Notation}
%\theoremstyle{remark}
\newtheorem{Remark}[Theorem]{Remark}
\newcommand{\hankredu}{\mathcal{H}_{ \alpha, \beta }}
\newcommand{\QNUM}{D}


\usepackage{subcaption}
\usepackage{float}
%\usepackage{xcolor}
\usepackage{bm}
\usepackage{hyperref}


\newcommand{\norm}[1]{\left\lVert#1\right\rVert}
\newcommand{\abs}[1]{\left|#1\right|}

\newcommand{\x}{\mathbf{x}}
\newcommand{\y}{y}
\newcommand{\ysh}{y_{\Sigma}}
\newcommand{\bb}{\mathbf{b}}
\newcommand{\p}{\mathbf{p}}
%\newcommand{\pu}{\mathbf{p}^{\mathbf{u}}}
\newcommand{\pu}{\mathbf{p}}
\newcommand{\uu}{\mathbf{u}}
\newcommand{\fiu}{\mathbf{\Phi}^{\mathbf{u}}}
\newcommand{\w}{w}
\newcommand{\e}{e}
\newcommand{\s}{\mathbf{s}}
\newcommand{\bt}{\boldsymbol{\tau}}
\newcommand{\Mk}{M_{w_1,\ldots, w_k}}
\newcommand{\hMk}{\widehat{\mathbf{M}}_{w_1,\ldots,w_k}}
\newcommand{\Mw}{M_{w}}
\newcommand{\hMw}{\widehat{\mathbf{M}}_{w}}

\newcommand{\htMw}{\widehat{\widetilde{\mathbf{M}}}_w}
\newcommand{\tMw}{\widetilde{M}_w}

\newcommand{\bu}{\mathbf{u}}
\newcommand{\bq}{\mathbf{q}}
\newcommand{\bx}{\mathbf{x}}
\newcommand{\bX}{\mathbf{X}}
\newcommand{\by}{\mathbf{y}}
\newcommand{\bz}{\mathbf{z}}
\newcommand{\bw}{\mathbf{w}}
\newcommand{\bv}{\mathbf{v}}
\newcommand{\bchi}{\bm{\chi}}
\newcommand{\bxi}{\bm{\xi}}
\newcommand{\bXi}{\bm{\Xi}}
\newcommand{\bHhat}{\mathbf{\widehat{H}}}
\newcommand{\bAhat}{\mathbf{\widehat{A}}}
\newcommand{\bBhat}{\mathbf{\widehat{B}}}
\newcommand{\bChat}{\mathbf{\widehat{C}}}
\newcommand{\bZhat}{\mathbf{\widehat{Z}}}
\newcommand{\sat}{\mathrm{sat}}
\newcommand{\SPAN}{\mathrm{Span}}
\newcommand{\Rank}{\mathrm{Rank}}

\pagestyle{empty}
\begin{document}


\def\BibTeX{{\rm B\kern-.05em{\sc i\kern-.025em b}\kern-.08em
    T\kern-.1667em\lower.7ex\hbox{E}\kern-.125emX}}
\markboth{\journalname, VOL. XX, NO. XX, XXXX 2017}
{Author \MakeLowercase{M. Petreczky \textit{et al.}}}

\title{Stability of input-output maps and their minimal realizations  in  state-linear, state-affine, LPV, and linear switched systems}

\author{Mihály Petreczky, Juan-Pablo Ortega,
Florian Rossmannek and 
Bálint Daróczy
%\thanks{$^{1}$Univ. Lille, CNRS, Centrale Lille, UMR 9189 CRIStAL, F-59000 Lille, France}
%\thanks{$^{2}$SPMS, NTU, Singapore}
%\thanks{$^{3}$HUN-REN SZTAKI, Budapest, Hungary}}
\thanks{M. Petreczky is with Univ. Lille, CNRS,
 Centrale Lille, UMR 9189 CRIStAL, F-59000 Lille, France
  (mihaly.petreczky@centralelille.fr).}
\thanks{J.-P. Ortega and 
F. \textcolor{black}{Rossmannek} are with SPMS, NTU, Singapore
  (e-mail: juan-{pablo.ortega,florian.\textcolor{black}{rossmannek}}@ntu.edu.sg).}  
%\thanks{F. Rossmannek is with SPMS, NTU, Singapore
%  (e-mail: florian.rossmannek@ntu.edu.sg).}  
\thanks{B. Daróczy is with
HUN-REN SZTAKI, Budapest, Hungary (e-mail: daroczy.balint@sztaki.hun-ren.hu).}}

%\institute[]{\texorpdfstring{$^1$}{1}University of Lille, CNRS,
%Centrale Lille, UMR 9189 CRIStAL, F-59000 Lille, France\\
%\texorpdfstring{$^3$}{3}Division of Mathematical Sciences, School of Physical and Math-
%ematical Sciences, Nanyang Technological University, Singapore
%\\
%\texorpdfstring{$^3$}{3}HUN-REN Institute for Computer Science and
%Control HUN-REN SZTAKI, Budapest, Hungary}

%\begin{document}
\maketitle 
\thispagestyle{empty}

\begin{abstract}
      Stability is often assumed in learning and identification, yet it is rarely characterized directly from input--output data. We show that an input--output family admits a stable finite-dimensional state-linear realization iff it has finite Hankel-rank and its response decays uniformly with time; for state-linear realizable maps this decay is necessarily exponential. We extend these results to state-affine, LPV, and linear switched systems via suitable input-forgetting notions, and relate forgetting to decay of impulse responses (sub-Markov parameters). In all cases, the decay/forgetting rate determines the decay rate of every minimal realization.
\end{abstract}
\begin{IEEEkeywords}
Stability, input forgetting, minimal realization, state-affine systems, LPV systems. 
\end{IEEEkeywords}
\section{Introduction}
In this paper we are interested in characterizing input-output maps
\textcolor{black}{of stable \textcolor{black}{state-linear} systems (SLS)} %state-linear systems (SLS):
\begin{equation}
    \label{eq1} 
       \begin{split}
         &  x(t+1)=A_{u(t)}x(t), \quad x(0) \in B, \quad  y(t)=C x(t)
      \end{split}
\end{equation}
where $x(t) \in \mathbb{R}^n$ is the state, $y(t) \in \mathbb{R}^{\textcolor{black}{n_y}}$ is the output and $u(t) \in \mathbb{U}$ is the 
\textcolor{black}{input at time $t$}, $A_{u(t)}$ and $C$ are matrices of suitable sizes,
% \in \mathbb{R}^{n \times n}$, $b_{u(t)} \in \mathbb{R}^{n}$, 
   %$C \in \mathbb{R}^{p \times n}$ 
   and $B=\{B_j \in \mathbb{R}^n\}_{j \in J}$
   is a family of initial states indexed by a set $J$.
   %and $\{A_u\}_{u \in \mathbb{U}}$ is a bounded set. 
  \textcolor{black}{SLSs are also known as switched linear systems (with no additive input) when $\mathbb{U}$ is finite or $\{A_u\}_{u \in \mathbb{U}}$ is compact. We prefer the term SLS as it is older \cite{Son:Real} and it is not associated with 
  assumptions on $\mathbb{U}$.}
  %$\{A_u\}_{u \in \mathbb{U}}$.}%being finite or 
  %$\{A_u\}_{u \in \mathbb{U}}$ is compact.}
  % 
  % family of matrices indexed by $\mathbb{U}$, i.e., for every $u \in \mathbb{U}$ there exists a matrix $A_u$ such that $A_{u(t)}=A_{u}$ if $u(t)=u$.
   %µ Note that for different indices $j_1,j_2 \in J$, $B_{j_1}=B_{j_2}$ is allowed.
   % $x(t) \in \mathbb{R}^n$ is the state, $u(t) \in \mathbb{U} \subseteq \mathbb{R}^{m}$ is
   %the input and $y(t) \in \mathcal{Y} \subseteq \mathbb{R}^p$. 
   %\textcolor{black}{By stability of \eqref{eq1} we mean global uniform exponential stability to zero (\emph{GUES}).}


  \begin{color}{black}
   Stability of \eqref{eq1} is well-understood, without claiming completeness \cite{Chitour1,RJungers,IanMorris2022,Protasov2022,mason2023}. In particular, it is known that the asymptotic stability of \eqref{eq1} is equivalent to uniform \emph{global exponential stability (GUES)} and both are equivalent to the joint spectral radius of $\{A_u\}_{u \in \mathbb{U}}$ being less than one. 
   However, it is less clear how stability properties of the input-output behavior of \eqref{eq1} relate to the stability of \eqref{eq1}, e.g., it is unclear if
   \eqref{eq1} is GUES if its input-output maps decay.
   %if the input-output maps of \eqref{eq1} decaying, then 
   
   %of \eqref{eq1} relate to those of 
   %\eqref{eq1}. 
   %For instance, if \eqref{eq1} is GUES, then its input-output maps are exponentially decaying,
   %but it is  unclear if and when the converse holds.
   % i.e. that any SLS realization of a family of (exponentially) decaying input-output maps is necessarily asyptotically stable or GUES. 
   %Classical results for linear systems suggest that 
   %this should be the case for minimal realizations.
  Our \emph{first contribution} is to show that a family of input-output maps has a GUES SLS realization if and only if it has an SLS realization and the values of the input-output maps decay to zero; in that case, decay is necessarily
  exponential, and all minimal realizations are GUES with this decay rate.

  This is relevant for two reasons. 
  First, state-space representations (SSRs) obtained through identification or model reduction capture, at best, only the true input-output behavior, 
  and therefore cannot be identified with the true system itself. 
  Hence, in the absence of prior knowledge, any SSR reproducing the true input-output behavior is an equally valid candidate model.
  %First, state-space representations (SSRs) obtained via identification or model reduction encode at best only the true input-output behavior, 
  %and hence they cannot be identified with the true system. Rather, any SSR whose input-output behavior is that of the true system is a valid candidate model for the true system, unless prior knowledge is available.  
  Therefore, analysis and control design for one such SSR should apply to 
  other SSRs with the same input-output map, e.g.,
  a controller that stabilizes one SSR (and thus its input-output behavior) should stabilize other SSRs with the same input-output behavior. 
  Our results suggest that this might be the case if the other SSRs are minimal.
  Second, in learning/system identification, stability is often imposed on SSRs to guarantee robust long-term prediction and prevent error accumulation. However, this restriction is only meaningful if the true input-output behavior admits a stable realization. Our results provide conditions for this.
  %and ensure that all minimal realizations inherit the same stability properties.

  As a \emph{second contribution} we apply 
  %is the application of 
  the results for SLSs 
  %on equivalence between stability of input-output maps and their SLSs realizations 
  to \emph{state-affine systems (SAS) \cite{Son:Real}, linear parameter-varying systems with affine dependence on the scheduling variable (LPV-SSA) \cite{Petreczky2016,Toth1}, and linear switched systems (LSS) \cite{Sun:Book}}.
  Realizability by a GUES \emph{SAS }is equivalent to an appropriate input-forgetting property, where a SAS is called GUES if its linearization is GUES. 
  Realizability by a GUES LPV-SSA/LSS, i.e., one for which the SLS obtained by taking zero additive input is GUES, is equivalent to the impulse response (sub-Markov parameters) decaying to zero, and equivalently, to $\ell_\infty$-forgetting (input forgetting with BIBO stability).
  % the second contribution 
  %for \emph{state-affine (SAS), linear parameter-varying systems with affine dependence on the scheduling variable (LPV-SSA), and linear switched systems (LSS)}, the second contribution 
%
%SAS, LPV-SSA, and LSS reduce to the SLS case, analogous results hold
  %for those system classes 
   In all cases, the decay/forgetting rate is necessarily exponential and it coincides with the decay rate of minimal realizations, which are necessarily GUES. 
  %In particular,  our results
  %for LPV-SSA characterize completely those input-output maps which can be realized by exponentially stable LPV-SSAs, which
  %s useful for system identification for LPV-
  The relevance of these results for SASs/LPV-SSAs/LSSs mirrors that for SLSs. For SASs, they
  are also relevant for reservoir computing \cite{Ortega4,Ortega3,Sepulchre1,GONON202110}, as they relate
  the \textcolor{black}{fading memory and stability/echo-state properties.}
  % which is important in reservoir computing
  %A related motivation comes from reservoir computing, where stability-type assumptions ensure well-posedness for infinite input sequences and are tied to . Our results might help to make these connections explicit.
  \end{color}
%Informally, we call \eqref{gen:eq1} \emph{globally uniformly exponentially stable (GUES)}, if 
%any solution of \eqref{eq1} converges to zero with an exponential rate.
%there exists $K > 0$, $\beta \in (0,1)$, such that
%if for any solution $x$ of  \eqref{eq1}, including those for which $x(0) \notin B$,
%the dynamical system %$x(t+1)=f(x(t),u(t))$
%   \begin{equation}
%    \label{eq2}
%     x(t+1)=A_{u(t)}x(t)
 %  \end{equation} 
  %is uniformly exponentially stable at zero for any choice of $u(t)$, i.e., 
  %if there exists $K > 0$, $\beta \in (0,1)$, such that 
  %$\|x(t)\| \le K\beta^{t} \|x(0)\|$.
  % for any solution of (including those for which $x(0) \notin B$).

%  \textbf{Contributions.}
%µ  \textbf{(A) SLS:} 
  %\textbf{((B) Decay $\Rightarrow$ exponential:} For SLS-realizable maps, uniform decay is equivalent to \emph{uniform exponential} decay.\\
  %\textbf{Minimal realizations inherit stability:} If the response decays exponentially with rate $\beta$, then \emph{every minimal} SLS realization is GUES with the \emph{same} rate $\beta$. \\
  %\textcolor{black}{\textbf{(B) Extensions to state-affine (SAS), linear parameter-varying systems with affine dependence on the scheduling variable (LPV-SSA), and linear switched systems (LSS):}
  
  %\\
  %\textbf{LPV/LSS impulse-response view:} GUES realizability is equivalent to finite Hankel-rank and \emph{uniform decay of sub-Markov parameters} (impulse responses), and the decay rate 
  %\textbf{(B.3) } For LPV/LSS we relate decay of impulse response  to \emph{$\ell_\infty$-forgetting}, which combines 

%  \begin{color}{black}
%  From a technical point of view, the main contribution is a non-trivial combination of realization theory of 
%    SLSs and LPV/LSS systems and the characterization of stability in terms of joint spectral radius of the system matrices. 
%  \end{color}
  %\textbf{Motivation and significance.}
%%  \begin{color}{black}
%%  On a conceptual level, if an SSR results from system identification or model reduction, it represents only the input-output behavior of the true system. Thus, analysis/control design should ideally apply to all SSRs with the same input-output behavior. For example, if feedback designed using one SSR stabilizes this SSR and hence the input-output behavior, we should know whether it stabilizes other SSRs which are input-output equivalent. 
%%
%%  On a practical level, learning from data involves finding the best SSR from a parametrized family that fits the training data. For technical reasons, stability is often assumed: small input perturbations should cause only small output changes over longer horizons. This ensures robust prediction based on past inputs and avoids numerical error accumulation. However, if we restrict to stable SSRs, we must ensure the true input-output behavior admits such a realization. Without such knowledge, we cannot guarantee that requiring exponentially stable SSRs won't exclude an exact realization of the input-output behavior. Our results clarify this correspondence: they characterize when input-output behavior can be realized by stable systems and ensure that minimal realizations inherit the same stability properties.
%%
%%  A related motivation comes from reservoir computing, where stability-type assumptions ensure well-posedness for infinite input sequences and are tied to fading memory/echo-state properties \cite{Ortega1,Ortega2,Sepulchre1,GONON202110}. Our results might help to make these connections explicit.
%%  \end{color}
  
  %At the same time, preservation of various notions of stability is a fundamental requirement for learning and system identification algorithms, and it is important to understand how the stability properties of the system are related to the properties of the response map.
  %In this paper we will be interested in characterizing input-output maps which 
  %can be realized by a stable state-affine system. First, we examine the input-output maps of \emph{state-linear systems (SLS)}, a special case of SSR, where the non zero initial states are not the broadest. Later on we characterize other subclasses of SSR.  
%   
   %\begin{equation}
   % \label{eq1} 
   %    \begin{split}
   %      &  x(t+1)=A_{u(t)}x(t), \quad x(0) \in B\\
    %     & y(t)=C x(t)
    %   \end{split}
   %\end{equation}
   %where $A_{u(t)} \in \mathbb{R}^{n \times n}$, $b_{u(t)} \in \mathbb{R}^{n}$, 
   %$C \in \mathbb{R}^{p \times n}$, $d_{u(t)}$, 
   % $x(t) \in \mathbb{R}^n$ is the state, $u(t) \in \mathbb{U} \subseteq \mathbb{R}^{m}$ is
   %the input and $y(t) \in \mathcal{Y} \subseteq \mathbb{R}^p$, 
   %and $B$ is a set of initial states.
%, i.e., 
   %In this paper we consider input-output maps of \emph{state-linear systems (SLS)}
   % In this paper we consider input-output maps of \emph{state-linear systems (SLS)}
   % Moreover,
   %we will assume that there exists functions $\psi_i$, $i=1,\ldots,n_p$ such that
   
   %$F:\mathbb{U} \rightarrow \mathbb{R}^{n \times n}$, $b:\mathbb{U} \righatrrow \mathbb{R}^n$
   
   %are polynomial, i.e. for every $u \in \mathbb{U}$.
   %For the sake of simplicity, we consider 
%
\textbf{Related work.}
  For linear systems, the relationships among internal stability, BIBO stability, and $\ell_p$-gains are classical results \cite{Cal:Des}. For nonlinear systems, similar connections are subtler and often require strong reachability/observability assumptions \cite{WillemsStabOld,FROMION1996243,van19922,HILL1980327,MoylanCounterExamples,478341}, which may not hold for structured classes like SLS/SAS/LPV-SSA/LSS.
  \begin{color}{black}
  %We do not claim novelty with respect to the literature on stability of SLS/SAS/LPV-SSA/LSS state-space representations
  While stability of SLS/SAS/LPV-SSA/LSS state-space representations is well-studied, e.g., \cite{Chitour1,RJungers,IanMorris2022,Protasov2022,mason2023,Toth1,Sun:Book},
  %which, in particular, establish the equivalence of asymptotic and exponential stability and the role of joint spectral radius in stability. However, these works do not fully characterize when an input--output map admits a stable realization. . 
  our contribution is to establish the equivalence between decay/input-forgetting of input-output maps
  and stability properties of their minimal realizations, %which is shown
  %necessary and sufficient input--output conditions for stable realizability and to show that minimal realizations inherit stability properties, 
  by combining realization theory with known stability results for SSRs.
  \end{color}
  In reservoir computing, 
  fading memory and stability 
  were studied for SASs \cite{Ortega4,Ortega3,Sepulchre1,GONON202110}, 
  \textcolor{black}{but not the correspondence between 
  fading-memory of
  input-output maps and the stability of their
  minimal finite-dimensional realizations. }
  %in the realization-theoretic framework presented here.
  
  



%\begin{Notation}
 %\label{label:not} 

\section{Preliminaries}
 A \emph{nonempty word} over $X$ is a finite sequence of elements of $X$, i.e.,  
$w = x_{1}x_{2}\cdots x_{k}$ for some $k \in \mathbb{N}$, $k >0$, $x_{1}, x_{2}, \ldots, x_{k} \in X$; $|w|:=k$ is the 
length of $w$. \textcolor{black}{The set of all nonempty words is $X^{+}$. The empty word is denoted by $\epsilon$, $|\epsilon|=0$, and $X^{*} = \{\epsilon\} \cup X^{+}$. Concatenation is defined by $vw = a_1\cdots a_m b_1\cdots b_n$ for $v = a_1\cdots a_m$, $w = b_1\cdots b_n$, and $v\epsilon = \epsilon v = v$.}
We denote by $X^k$ the set of sequences of length $k$, and by $X^{\infty}$ the set of all infinite sequences of elements of $X$;
for any $w \in X^{\infty}$, $|w|=\infty$.
If $w \in X^{*} \cup X^{\infty}$, and $k \le |w|$ then $w[k]$ denotes the $k$th element of 
$w$, i.e., $w=w[1]\cdots w[|w|]$. \textcolor{black}{If $i > |w|$, then $w[i]\cdots w[|w|]$ is understood to be the empty word $\epsilon$.} 
	\textcolor{black}{If $Z$ is an $n \times m$ matrix, then $Z_{i,j}$ denotes the element on the $i$th row and $j$th column of $Z$; $Z_{i,\cdot}$ and $Z_{\cdot,j}$ denote the $i$th row and $j$th column of $Z$; if $z$ is a vector \textcolor{red}{with $n$ entries}, then $z_i$ denotes the $i$th element of $z$. 
%\textcolor{black}{if $w$ is a sequence of vectors, then $w_i[k]$ is the $i$th entry of the vector $w[k]$.}
We denote by $\|Z\|_F$ the Frobenius norm of $Z$, by $\|Z\|_2$ the spectral norm of $Z$,  \textcolor{red}{by $\|z\|_2$ the Euclidean norm of $z$, and 
by $\|z\|_{\infty}$ the supremum norm, i.e., $\|z\|_{\infty}=\max_{i,\ldots,n} |z_i|$.}  Let $I_n$ be the $n \times n$ identity matrix.}
\textcolor{black}{For a sequence $w \in (\mathbb{R}^p)^{*} \cup (\mathbb{R}^p)^{\infty}$ 
%such that $X$ is Euclidean space $\mathbb{R}^n$,
define  $\|w\|_{\infty} = \sup_{1 \le i \le |w|, i \in \mathbb{N}} \|w[i]\|_{\infty}$.}
For a  family $\{A_x\}_{x \in X}$ of $n \times n$ matrices 
and a sequence
$w  \in X^*$ define 
\[ A_{w} = A_{w[|w|]} \cdots A_{w[1]}, ~~\textcolor{black}{\mbox{ if } w \in X^{+}}, \quad A_{\epsilon}=I_n. \]
%and let $A_{\epsilon}=I_n$ be the identity matrix. 

%\end{Notation}

\subsection{State-space representations}
   %We identify a LSL of the form \eqref{eq1} with the tuple
   %$\Sigma=(n,\{A_u\}_{u \in \mathbb{U}},B,C)$. 
   %In addition, we will assume that 
   Consider a general state-space representation (\emph{SSR})
\begin{equation}
\label{gen:eq1}
\Sigma \left \{\begin{split}
 x(t+1)&=f(x(t),u(t)), x(0) \in B \\
 y(t)& =h(x(t),u(t))
\end{split}\right. \textcolor{black}{,}
\end{equation}
where $x(t) \in \mathbb{R}^n$ is the state, $u(t) \in \mathbb{U}$ is the input, $y(t) \in \mathbb{R}^{\textcolor{black}{n_y}}$ is the output,
$f:\mathbb{R}^n \times \mathbb{U} \rightarrow \mathbb{R}^n$, 
$h:\mathbb{R}^n  \times \mathbb{U} \rightarrow \mathbb{R}^{\textcolor{black}{n_y}}$ is the \emph{state-transition} and
\emph{readout} map respectively, and  
   $B=\{B_j \in \mathbb{R}^n \}_{j \in J}$, \textcolor{black}{$J \ne \emptyset$}, is the family of \emph{initial states}. 
We identify  $\Sigma$ with the tuple $(f,B,h)$ and we denote by  $\dim(\Sigma)$ the dimension $n$ of its state-space.
%We denote the dimension of $\Sigma$ by $n=\dim(\Sigma)$.
%In order to avoid 
%   excessive notation, we assume there exists an index set $J$ and a map $\mu:J \rightarrow B$ and will identify 
%   
%   $B$ with the family of indexed elements
%   Note that for different indices $j_1,j_2 \in J$, $B_{j_1}=B_{j_2}$ is allowed. 
%We will refer to \eqref{gen:eq1} in this case as a $J$-SSR. Most of the time $J$
%will be clear from the context, and we will just use the term SSR.
%We call $n$ the dimension of  by $\dim(\Sigma)$.

For an initial state $x_0$,  and sequence $w \in \mathbb{U}^{*}$
denote by $x_{\Sigma}(x_0,w)$ the state of \eqref{gen:eq1} 
reached at time $t$ under input $u(0)=w[1],\ldots, u(t-1)=w[t]$, $t=|w|$ from \textcolor{red}{the} initial state $x_0$, and we set 
$x_{\Sigma}(x_0,\epsilon)=x_0$.
Define the response map $y_{\Sigma,x_0}\textcolor{black}{:\mathbb{U}^{+} \rightarrow \mathbb{R}^{n_y}}$ of $\Sigma$ induced by the
initial state $x_0$ as follows: \textcolor{black}{$y_{\Sigma,x_0}(v\sigma)=h(x_{\Sigma}(x_0,v),\sigma)$,
$v \in \mathbb{U}^{*},\sigma \in \mathbb{U}$, i.e.,
  $y_{\Sigma,x_0}(u(0)\cdots u(t))$ is the output response of
  \eqref{gen:eq1} at time $t$ to the inputs $\{u(s)\}_{s=0}^{t}$ from the initial state $x_0$.
}
%$y_{\Sigma,x_0}: \mathbb{U}^{+} \ni w \mapsto h(x_{\Sigma}(x_0,\textcolor{black}{w[1]\cdots w[|w|-1]}),w[|w|])$, i.e,
%$y_{\Sigma,x_0}(w)$ is the output $y(t)$ under inputs \textcolor{black}{$u(0)=w[1],\ldots, u(t)=w[t+1]$, $|w|=t+1$} from initial state $x_0$.
%Among possible initial states we will be interested in the zero initial states. 
%Following \cite{Son:Real}, 
We say that $\Sigma$ is \emph{span-reachable} \textcolor{black}{(respectively \emph{affine-reachable})}, 
if the linear (respectively affine) span of all the reachable states $\{x_{\Sigma}(B_j,w)\}_{w \in \mathbb{U}^{*}, B_j \in B}$
equals $\mathbb{R}^n$. We say that $\Sigma$ is \emph{observable}, if 
for any two distinct initial states $x_0 \ne \bar{x}_0$ there exists $w \in \mathbb{U}^{+}$ such that
$y_{\Sigma,x_0}(w) \ne y_{\Sigma,\bar{x}_0}(w)$. 

%We consider several subclasses of SSR.


\emph{State-linear systems (SLSs)} are special cases of SSR \eqref{gen:eq1} 
for which $f(x,u)=A_u x$ and $h(x,u)=Cx$. In what follows, we identify $\Sigma$ with the tuple 
$\Sigma=(\{A_u\}_{u \in \mathbb{U}},B,C)$. 
   
\emph{State-affine systems (SAS) \cite{Son:Real}} are SSRs \eqref{gen:eq1} of the form
\begin{equation}
\label{eq1:sas}
%\Sigma \left \{
\begin{split}
 x(t+1)&= A_{u(t)}x(t)+b_{u(t)}, \quad  x(0) \in B=\{x_0\}, \\
 y(t) &= Cx(t)+d,
\end{split}
%\right.
\end{equation}
where $A_u \in \mathbb{R}^{n \times n}$, $b_u \in \mathbb{R}^{n}$ for $u \in \mathbb{U}$, and
$C \in \mathbb{R}^{\textcolor{black}{n_y} \times n}$, $d \in \mathbb{R}^{\textcolor{black}{n_y}}$, and the family of initial states $B$ is a singleton.
% then we speak of a \emph{state-affine system (SAS)}.
%In the sequel, we will also be interested in the following subsets of SASs and response maps,
Following \cite{Son:Real} we define the subclass of bounded  SASs, for 
which realization theory is computationally effective.  To this end, 
let $0 < n_p \in \mathbb{N}$ %and $Q=\{0,\ldots,n_p\}$ 
and consider $n_p+1$
linearly independent  functions $\psi_{q}:\mathbb{U} \rightarrow \mathbb{R}$, 
$q=0,\ldots,n_p$. % We denote by $\psi$ the collection $\{\psi_q\}_{q \in Q}$.
Then  SSR of \eqref{eq1:sas} is \emph{$\{\psi_q\}_{q=0}^{n_p}$-bounded SAS}, if 
%$B=\{0\}$ and 
%\begin{equation}
%µ\label{eq:sas_psi}
%\begin{split}
$A_ux+b_u\!\!=\!\!\sum_{i=0}^{n_p} (A_ix+b_i) \psi_i(u)$ and  $h(x,u)\!\!=\!\!Cx+d$
%\\ 
%~ h(x,u)=\sum_{i=0}^{n_p} (C_ix+d_i)\psi_i(u(t),  \\
%\quad d_{u(t)}=\sum_{i=1}^{n_p} d_i \psi_i(u(t)) \\
%\end{split}
%\end{equation} 
for \textcolor{black}{some} matrices and vectors 
$\{A_q,b_q\}_{q=0}^{n_p}$.
\emph{Bilinear systems \cite{isi:tac}} are bounded SASs 
\textcolor{black}{with $\mathbb{U}=\mathbb{R}^{m}$, $\psi_0=1$, $\psi_q(u)=u_q$ for $q \in \{1,\ldots,m\}$, $m=n_p$, $b_0=0$, $d_0=0$.}
%$q=1,\ldots, m$ for all $u=(u_1,\ldots,u_m)^T$, and 
%and $\psi_q(u)$ is the $q$th entry of $u$, and %$q=1,\ldots, m$ for all $u=(u_1,\ldots,u_m)^T$, and 
%$b_0=0$, $d=0$ i.e.,  $A_u=A_0+\sum_{i=1}^{m} A_iu_i$ and $b_u=\sum_{i=1}^{m} b_iu_i$.
% $d_0=0$, $C_i=0,d_i=0$, $i=1,\ldots,n_p$,
%$h(x,u)$ is linear and depends only on $x$, 
%bounded SASs are called \emph{bilinear systems}. In this case, 
 %$f(x,u)=A_0x + \sum_{i=1}^{m} u_i A_ix + \sum_{i=1}^{m} u_i b_i$, $h(x,u)=Cx$ for suitable matrices. 

 %If $\mathbb{U}=Q \times \mathbb{R}^m$ and $Q=\{1,\ldots, \QNUM\}$ is a finite set, 
%and for $u(t)=(q(t),v(t))$, $q(t)\in Q$, $v(t)\in\mathbb{R}^m$, the SSR \eqref{gen:eq1} takes the form
%
%for suitable matrices $A_q,B_q,C_q$ for all $q\in Q$, then
% we call \eqref{gen:eq1} a \emph{linear switched system (LSS)}. 
 
 A \emph{linear parameter-varying (LPV) system with affine dependence on the parameters (LPV-SSA) \cite{Toth1}} is an SSR
 %are special cases of SASs, where the scheduling signal $p$ is viewed as the input signal, and the control input $v$ is constant $1$. In particular,
 with inputs $u(t)=(p(t),v(t)) \in \mathbb{U}=\mathbb{P} \times \mathbb{R}^m$ of the form
 %and the state update and 
 %and with $A(p(t))
 % where $\mathbb{P}$ is a bounded
 %subset of $\mathbb{R}^{n_p}$ %(say $\mathbb{P}=[-1,1]^{n_{\mathbb{p}}})$ and 
\begin{equation}
\label{eq1:lpv}
%\Sigma : \begin{cases}
\begin{split}
x(t+1) & = A(p(t)) x(t) + B(p(t)) v(t), x(0)=0\\
%(A_0 + \sum_{i=1}^{n_p} A_i p_i(t)) x(t) + (B_0 + \sum_{i=1}^{n_p} B_i p_i(t)) u(t), x(0)=0\\
y(t) &= C(p(t))x(t) \textcolor{black}{,}
%(C_0 + \sum_{i=1}^{n_p} C_i p_i(t)) x(t) %+ d(p(t)) u(t)
%\end{cases}
\end{split}
\end{equation}
where $p$ is called the \emph{scheduling signal}, $v$ is the control input, and 
the matrices $A(p(t))$, $B(p(t))$, and $C(p(t))$ are affine in $p(t)$, i.e.
\(
   X(p(t)) = X_0 + \sum_{i=1}^{n_p} X_i p_i(t)
\)
with $X = A, B, C$ and $\{A_i, B_i, C_i\}_{i=1}^{n_p}$ being suitable matrices,
\textcolor{black}{and the family of initial states is the singleton $\{0\}$}.
We assume that $\mathbb{P} \subseteq \mathbb{R}^{n_p}$ is contained in the standard simplex $\Delta = \{x = (x_1, \ldots, x_{n_p})^T \in [0,1]^{n_p} \mid \sum_i x_i = 1\}$
and $\mathbb{P}$ contains the standard unit vectors $e_1, \ldots, e_{n_p}$.
%This assumption can be achieved for any $\mathbb{P}$ which does not lie in proper subspace and 
%which is contained in a polytope,
%if reparametrize it by barycentric coordinates, so that
%$X(p)=\sum_{i=1}^{K} p_i X_i$ for $p\in\Delta$, where $X_i$ are the vertex values of $X\in\{A,B,C\}$.
We also assume $A_0 = 0$, $B_0 = 0$, $C_0 = 0$, otherwise we can replace $A_i, B_i, C_i$ by $A_i + A_0, B_i + B_0, C_i + C_0$ for $i = 1, \ldots, n_p$.
A \emph{linear switched system (LSS) \cite{Sun:Book,D:Lib}} is a special case of LPV-SSA where $\mathbb{P} = \{e_1, \ldots, e_{n_p}\}$,
 %where $e_i$ is the $i$-th standard unit vector in $\mathbb{R}^{n_p}$.
 \textcolor{black}{and  $X(e_i) = X_i$ for $X = A, B, C$ and $i = 1, \ldots, n_p$, and the scheduling signal $p$ is identified with the
  \emph{switching signal} $q$, where $p(t) = e_{q(t)}$ for all $t \in \mathbb{N}$.  }
%i.e., \eqref{eq1:lpv} becomes:
%\eqref{eq1:lss}.
%\begin{equation}
%\label{eq:lss}
%\begin{split}
%x(t+1) &= A_{q(t)}x(t)+B_{q(t)}v(t), \quad x(0)=0,\\
% y(t) &= C_{q(t)}x(t).
%\end{split}%\right.
%\end{equation}
%$x(t+1) = A_{q(t)} x(t) + B_{q(t)} v(t)$, $y(t) = C_{q(t)} x(t)$.    
%A \emph{bounded SAS} is a special case of an LPV where $\mathbb{P}=\{ (\psi_1(u),\ldots, \psi_{n_p}(u)) \mid u \in \mathbb{U}\}$, $m=1$, and instead of input sequences 
%$u=u(0)u(1)\cdots u(t-1)$, we consider scheduling sequences $p=p(0)p(1)\cdots p(t-1)$, where $p(t)=(\psi_1(u(t)),\ldots, \psi_{n_p}(u(t)))$ for all $t \in \mathbb{N}$. In this case, the system matrices are given by
%$A(p) = \sum_{i=1}^{n_p} F_i p_i$, $B(p) = \sum_{i=1}^{n_p} b_i p_i$, $C(p) = \sum_{i=1}^{n_p} C_i p_i$ for suitable matrices $F_i, b_i, C_i$, $i=1,\ldots,n_p$, and we consider only constant $1$ inputs. In particular, a bilinear system arises as a special case of an LPV system
%where the scheduling signal is viewed as the input signal, and the control input $v$ is constant $1$. 
%In fact, if $\mathbb{P}$ is a convex polytope generated by $K$ vertices $v_1, \ldots, v_K$, then it can always be identified with the standard simplex in $\mathbb{R}^K$ by identifying every element $p$ of $\mathbb{P}$ with the coordinates $(\alpha_1, \ldots, \alpha_K)$ of the convex combination $p = \sum_{i=1}^{K} \alpha_i v_i$. In the new space of scheduling variables, $v_i$ become the unit vectors, and the system matrices can be taken as $X(\alpha) = \sum_{i=1}^{K} \alpha_i X_i(v_i)$ for $X = A, B, C$. That is, $\{A_i, B_i, C_i\}_{i=1}^{K}$ are the values of the original matrices $A, B, C$ at the vertices of $\mathbb{P}$.
% Short replacement for the previous paragraph (you can delete the longer version above if desired).
%\begin{Remark}
%\end{Remark}
%\end{Remark}
%\begin{Remark}
%In fact, if $\mathbb{P}$ is a convex polytope generated by $K$ vertices $v_1, \ldots, v_K$, then it can always be identified with the standard simplex in $\mathbb{R}^K$ by identifying every element $p$ of $\mathbb{P}$ with the coordinates $(\alpha_1, \ldots, \alpha_K)$ of the convex combination $p = \sum_{i=1}^{K} \alpha_i v_i$. In the new space of scheduling variables, $v_i$ become the unit vectors, and the system matrices can be taken as $X(\alpha) = \sum_{i=1}^{K} \alpha_i X_i(v_i)$ for $X = A, B, C$. That is, $\{A_i, B_i, C_i\}_{i=1}^{K}$ are the values of the original matrices $A, B, C$ at the vertices of $\mathbb{P}$.
%\end{Remark}
 %$f(x,(p,v))=A_px+B_pv$, $h(x,(p,v))=C_px$  for suitable matrices $A_p,B_p,C_p$ 
 %for all $q \in Q$, $v \in \mathbb{R}^m$, then 
 %\eqref{gen:eq1} correspond to \emph{linear parameter-varying state-space representation (LPV-SS).}
 %In addition, if there exist matrices 
 %functions $A_i,B_i,C_i$, $i=0,\ldots,n_p$ such that
 %$\psi_1,\ldots,\psi_{n_p}:\mathbb{P} \rightarrow \mathbb{R}$ such that with 
 %$\psi_0,\psi_1,\ldots, \psi_{n_p}$ are linearly independent, where $\psi_0=1$,
 %and matrices $A_i,B_i,C_i$, $i=0,\ldots,n_p$ such that
 %$A(p)=A_0+\sum_{i=0}^{n_p} A_i p_i$, $B(p)=B_0+\sum_{i=0}^{n_p} B_i p_i$,
 %$C(p)=C_0+\sum_{i=0}^{n_p} C_i p_i$, then 
 %\eqref{gen:eq1} is a LPV 
 %\emph{state-space representation with affine dependence on the parameter (LPV-SSA)}. 
 %Without loss of generality, one can replace $\mathbb{P}$ by $\bar{P}=\{ (\psi_1(p),\ldots, \psi(p) \mid p \in \mathbb{P})\}$ and set $\psi_i(p)$ to be the $i$th coordinate of $p$. 

%For the cases of SAS, LSS, LPV-SSA we consider $B=\{0\}$,  i.e., there is a single initial state which is zero. 
%Let $\Sigma=(n,\{ A_{u} \}_{u \in \mathbb{U}},B,C)$,
%        \( \widetilde{\Sigma}=(\widetilde{n},\{ \widetilde{A}_{u}\}_{u \in \mathbb{U}}), \widetilde{B},\widetilde{C}\) be two LSL.
%Let $\widetilde{\Sigma}=$ be an SSR. 
Assume that both \textcolor{black}{$\widetilde{\Sigma}=(\widetilde{f},\widetilde{B},\widetilde{h})$} and $\Sigma$ belong to one of the classes: SLS, SASs,LPV-SSA, and LSS. Then $\widetilde{\Sigma}$ and $\Sigma$ \textcolor{black}{are called} \emph{isomorphic} if 
\textcolor{black}{$n\!\!=\!\!\dim(\Sigma)\!\!=\!\!\dim(\widetilde{\Sigma})$} and there exists an invertible map $T$ on $\mathbb{R}^n$ 
which is linear for SLSs, LSSs, LPV-SSAs,  and affine for SASs, such that
%and which maps the initial states of $\Sigma$ to the initial states of $\widetilde{\Sigma}$, i.e., $T(B_j)=\widetilde{B}_j$ for all $j \in J$, and
        %called a \emph{linear SSR isomorphism} from $\Sigma$ to 
        %$\widetilde{\Sigma}$, if %and is denoted by $T:\Sigma \rightarrow \widetilde{\Sigma}$
            %if $d_u=\widetilde{d}_u$, $u \in \mathbb{U}$, and 
            %which commutes with $f$, $B_j$, $h$, i.e.,
            %for all $u \in \mathbb{U}$.
            %the following equalities hold
            \begin{equation*}
                 \begin{array}{rcl}
                 \label{repr:morph:eq2}
                 &Tf(x,u)=\widetilde{f}(Tx,u), ~ TB_j=\widetilde{B}_j,  ~h(x,u)=\widetilde{h}(Tx,u) \textcolor{black}{,}
             \end{array}
            \end{equation*}
            \textcolor{black}{for all $x \in \mathbb{R}^n$, $u \in \mathbb{U}$ and $j \in J$.}
            %The SSR morphism $T$ is called \emph{surjective, injective, isomorphism}  
            %if $T$ is a surjective, injective or isomorphism respectively if viewed as a linear map.
            %µLinear SSR isorphisms between 
            %SLSs,  SASs,LSSs or LPV-SS(A)s will be called SLSs, SASs, LSSs, or 
            %LPV-SS(A) morphisms respectively. 
            
%\subsection{Response maps}
If the readout map $h$ in \eqref{gen:eq1} does not depend on $u$, then we call $\Sigma$ a \emph{state-output system (SO)}. In this case, the output at time $t$ depends only on the inputs up to $t-1$.
% i.e.,
%\begin{equation*}
%$y_{\Sigma,x_0}(wu)=h(x_{\Sigma}(x_0,w))$, for all $w\in\mathbb{U}^*,~u\in\mathbb{U},~x_0\in B$.
%\end{equation*}
We call functions $H:\mathbb{U}^{+}\to\mathbb{R}^{\textcolor{black}{n_y}}$ \emph{response maps}, and functions $H:\mathbb{U}^{*}\to\mathbb{R}^{\textcolor{black}{n_y}}$ \emph{SO response maps}. Intuitively, $H(u_0\cdots u_t)$ is the output at time $t$ (resp.\ at $t+1$ if $H$ is SO) generated by inputs $u_0,\ldots,u_t$, and for SO maps $H(\epsilon)$ is the output at time $0$. 
We model the input--output behavior of (SO) SSRs as the following families 
of (SO) response maps: % (resp.\ SO response maps).
% respectively \emph{SO response maps}.
\begin{align}
    \Psi=\{H_j:\mathbb{U}^{+}\rightarrow \mathbb{R}^{\textcolor{black}{n_y}}\}_{j \in J} & \mbox{ if general SSR} 
    \label{family:resp:gen}
    \\
    \Psi=\{ H_j:\mathbb{U}^{*}\rightarrow \mathbb{R}^{\textcolor{black}{n_y}}\}_{j \in J} & \mbox{ if SO SSR}  \textcolor{black}{.}
    \label{family:resp}
    %{H_j:\mathbb{U}^{+}\rightarrow \mathbb{R}^p\}_{j\in J} \quad \Psi=\{H_j:\mathbb{U}^{*}\rightarrow \mathbb{R}^p\}_{j\in J} \mbox{if SO}
\end{align}
The SSR $\Sigma$ is called a \emph{realization} of $\Psi$, if 
%$H_j(w)=y_{\Sigma,B_j}(w)$, and we say that $\Sigma$ is a realization of $\Psi$ from \eqref{family:resp} if $\Sigma$ is SO and for all $j \in J$, $w \in \mathbb{U}^*$, $H_j(w)=h(x_{\Sigma}(B_j,w))$.
\[
    H_j(w)=\left\{\begin{array}{rl}
             y_{\Sigma,B_j}(w) & \mbox{if $\Psi$ is from \eqref{family:resp:gen} and $w\in\mathbb{U}^{+}$} \\
	     h(x_{\Sigma}(B_j,w)) & \mbox{if $\Psi$ is from \eqref{family:resp} and $w \in \mathbb{U}^{*}$} \\
	                           & \textcolor{red}{\mbox{ and $\Sigma$ is SO}}
             \end{array}\right. 
\]
for all $j \in J$. %, $w \in \mathbb{U}^{*}$.
%
%from \eqref{family:resp:gen} if $H_j(w)=y_{\Sigma,B_j}(w)$ for all $w\in\mathbb{U}^{+}$, $j\in J$, and $\Sigma$ is a realization of $\Psi$ from \eqref{family:resp} if $\Sigma$ is SO and 
%$H_j(w)=h(x_{\Sigma}(B_j,w))$ for all $w\in\mathbb{U}^*$, $j\in J$. 
%and $H(\epsilon)$ correspond to the output of the underlying system at time $0$, and $H(u_0\cdots u_{t})$ represents the output response of  the system at time $t$ to
%or $H: \mathbb{U}^* \rightarrow \mathcal{Y}$,
%where $H(\epsilon)$ correspond to the output of the underlying system at time $0$, and
%\begin{align}
% 
%\end{align}  
%We model the input--output behavior of \eqref{gen:eq1} as a family
%\(
%\Psi=\{H_j:\mathbb{U}^{+}\rightarrow \mathcal{Y}\mid j\in J\}
%\)
%of \emph{response maps} \cite{Son:Real}, where $H_j(w)$ is the output at time
%$t=|w|$ generated by the input word $w\in\mathbb{U}^+$ from the initial state
%$B_j$. Formally, an SSR $\Sigma$ is a \emph{realization} of $\Psi$ if
%\(
%H_j(w)=y_{\Sigma,B_j}(w)
%\)
%for all $w\in\mathbb{U}^{+},~ j\in J$.
%Note that for SLSs and SASs the readout map $h(x,u)$ does not depend on $u$. Such SSR can realize only response maps for which the output is independent of the last input, i.e., $H_j(vq)=H_j(v)$ for all $v\in\mathbb{U}^*$, $q\in\mathbb{U}$, $j\in J$. In this case, we call $\phi:\mathbb{U}^+\to\mathcal{Y}$ \emph{state-output (SO)} if $\phi(vq)$ does not depend on $q$ for any $v\in\mathbb{U}^*$, and extend it to $\phi^{\mathrm{ext}}:\mathbb{U}^*\to\mathcal{Y}$ by $\phi^{\mathrm{ext}}(v)=\phi(vq)$ (well-defined by SO). We identify $H_j$ with $H_j^{\mathrm{ext}}$, and thus we work with maps on $\mathbb{U}^*$ and formalize input-output behaviors of SLSs as families of SO maps
%\begin{equation}
%   
%  \end{equation}
%  and an SSR realization of $\Psi$ is a realization of the family of SO maps in \eqref{family:resp}, 
%  if $H_j(w)=h(x_{\Sigma}(B_j,w))$ for all $w\in\mathbb{U}^*$, $j\in J$.
%We formalize potential  input-output behavior  of \eqref{gen:eq1} 
%as \emph{response maps} \cite{Son:Real}, i.e., functions of 
%of the form $H:\mathbb{U}^{+} \rightarrow \mathcal{Y}$, where %Intuitively, 
% %$H(\epsilon)$ correspond to the output of the underlying system at time $0$, and 
% $H(u_0\cdots u_{t})$ represents the output response of  the system at time $t$ to
% the inputs $u_0,\ldots, u_{t}$. We refer to such maps as \emph{response maps}. 
%Let $\Psi=\{H_j: \mathbb{U}^{+} \rightarrow \mathcal{Y} \mid j \in J\}$ be a 
%family of response maps indexed by $J$.
% The SSR \eqref{gen:eq1} is called a \emph{realization} of $\Psi$, if 
% $H_j(w)=y_{\Sigma}(B_j,w)$, i.e., $H_j(u(0) \cdots u(t-1))$
% is the output $y(t)$ of \eqref{gen:eq1} at time $t$ if the initial state is $B_j$ and 
% the inputs $u_0,\ldots, u_t$ were fed to \eqref{gen:eq1}. 
If $J$ is a singleton, we identify $\Psi=\{H\}$ with $H$ \textcolor{black}{and call} any SSR realization of $\Psi$  a realization of $H$.
Let $\mathscr{C}$ be a subclass of SSRs (e.g., SLS, SAS, LSS, LPV-SSA).
We call $\Sigma\in\mathscr{C}$ a \emph{minimal $\mathscr{C}$-realization} of $\Psi$ if \textcolor{black}{its dimension $\dim(\Sigma)$} is minimal %state-space dimension 
among all $\mathscr{C}$-realizations of $\Psi$.
 %Let $\mathscr{C}$ be a one of subclasses of SSRs, e.g., $\mathscr{C}$ is either the class of SLSs,SASs, LSS,LPV-SSs, LPV-SSAs. 
 %We call an  $\Sigma$ from $\mathscr{C}$ be  \emph{minimal ($\mathscr{C}$) realization} of $\Psi$, if it has the smallest state-space dimension
%among all  SSR realizations of $\Psi$ which belong to $\mathscr{C}$. 

%Next, we first deal only with SLSs and later on we adopt the result for other subclasses of SSR.

%In what follows we will first deal only with SSR of the form \eqref{eq1}, i.e., with SSRs which belong to the class of SLSs. Later, we will branch out to other classes such as SASs, LSSs,LPV-SSA(s). 
%

%In the rest of the paper we will assume that \eqref{eq1} is a $\psi$-SAS, i.e.,
%it satisfies \eqref{eq1.2}. We will then refere to $\psi$-SAS as SAS. 

%We then identify the SAS $\Sigma$ with the tuple $(n,\{F_q,b_q}_{q \in Q},H\})$.



%We say that \eqref{eq1} is \emph{polynomial}, if $F_u,b_u,C_u,d_u$ are polynomial functions
%of $u$.  
\subsection{Realization theory of state-linear systems}
\label{sect:real}
%Let $\Psi=\{H_j:\mathbb{U}^{+}\to\mathbb{Y}\mid j\in J\}$. If $\Psi$ is realized by the SLS \eqref{eq1}, then
%  $H_j(vq)=CA_vB_j$ for all $v\in\mathbb{U}^*$, $q\in\mathbb{U}$, $j\in J$; hence each $H_j$ must be independent of the last input.
%  We call $\phi:\mathbb{U}^+\to\mathbb{Y}$ \emph{state-output (SO)} if $\phi(vq)$ does not depend on $q$ for any $v\in\mathbb{U}^*$, and extend it to $\phi^{\mathrm{ext}}:\mathbb{U}^*\to\mathbb{Y}$ by $\phi^{\mathrm{ext}}(v)=\phi(vq)$ (well-defined by SO). We identify $H_j$ with $H_j^{\mathrm{ext}}$, and thus we work with maps on $\mathbb{U}^*$ and formalize input-output behaviors of SLSs as families:
  %Let $\Psi=\{ H_j:\mathbb{U}^{+} \rightarrow \mathbb{Y} \mid j \in J\}$. 
  %Note that if $\Psi$ is realized by a SLS \eqref{eq1}, then 
  %µ$H_j(vq)=CA_vB_j$ for all $v \in \mathbb{U}^{*}$, $q \in \mathbb{U}$, $j \in J$.
  %A necessary condition for $\Psi$ to be realizable by a SLS is that each $H_j\in\Psi$ is independent of the last input, i.e., $H_j(vq)$ does not depend on $q\in\mathbb{U}$ for any $v\in\mathbb{U}^*$.
  %We call $\phi:\mathbb{U}^+\to\mathbb{Y}$ \emph{state-output (SO)} if it has this property, and we extend it to $\phi^{\mathrm{ext}}:\mathbb{U}^*\to\mathbb{Y}$ by setting $\phi^{\mathrm{ext}}(v)=\phi(vq)$ for an arbitrary $q\in\mathbb{U}$ (well-defined by SO).
  %We identify $H_j$ with $H_j^{\mathrm{ext}}$, hence work with response maps on $\mathbb{U}^*$,
  Since the readout map of a SLS does not depend on the last input, we model the input-output behavior of SLSs as families of SO response maps
  \eqref{family:resp}. Then SLS \eqref{eq1} realizes $\Psi$, if and only if 
  \(
  H_j(w)=CA_wB_j,
  \), for all
  $w\in\mathbb{U}^*,~ j\in J$.
  We recall the necessary and 
  sufficient conditions for $\Psi$ being realizable by a SLS.
  To this end, define the shift of a response map $H:\mathbb{U}^{*} \rightarrow \mathbb{R}^{\textcolor{black}{n_y}}$ by sequence $v \in \mathbb{U}^{*}$ as follows:
  \[ v \circ H: \mathbb{U}^{*} \ni w \mapsto H(vw). \]
 We define the \emph{Hankel-space} $W_\Psi$ as the linear space generated by the shifted versions of elements of $\Psi$
 \[ 
     W_{\Psi}=\mathrm{Span} \{v \circ H_j \mid j \in J, v \in \mathbb{U}^{*} \} \textcolor{black}{,}
 \]
 and we call $\dim W_{\Psi}$ the \emph{Hankel-rank}. 
 \textcolor{black}{Here, for response maps addition and multiplication by a scalar
 is understood point-wise: for two SO response maps $H_1,H_2$, for scalars $\alpha,\beta \in \mathbb{R}$, $(\alpha H_1+\beta H_2)(w)=\alpha H_1(w)+\beta H_2(w)$, $w \in \mathbb{U}^{*}$.}
 %This terminology is motivated by the following observation:
 \textcolor{black}{If $\mathbb{U}$ and $J=\{j_1,\ldots,j_d\}$ 
 are both
 finite, then the}  SO response map $H_j$ corresponds to formal power series,
 and $W_{\Psi}$ is isomorphic to the column space of the Hankel-matrix
 \cite{MFliessHank,Reut:Book,Son:Real}
 of
 the formal power series $S:\mathbb{U}^{*} \ni v \mapsto \begin{bmatrix} H_{j_1}(v) & \ldots&  H_{j_d}(v) \end{bmatrix}$.
 %$H$ and $H$, i.e., 
% $\mathcal{W}_{H}$ is the linear space generated
% $\{H_v-H\}_{v \in \mathbb{U}^{*}}$.  We call $\mathcal{W}_H$ the \emph{shift-space of $H$}.
\begin{Theorem}[Theorem 2.15 \cite{Son:Real}]
\label{theo:real1}
      The family $\Psi$   
      has a realization by a SLS, if
      and only if $\dim W_{\Psi}=n <+\infty$.
      All minimal realizations of $\Psi$ are of dimension $n=\dim W_{\Psi}$, and 
      they are isomorphic.
      A SLS $\Sigma$ is a minimal realization of $\Psi$, if and only if 
      $\Sigma$ is span-reachable and observable. % and it is a realization of $\Psi$. 
      %All minimal SLS realizations of $\Psi$ are isomorphic. 
      %From any span-reachable SLS  realization $\Sigma_r$ of $\Psi$
      %and any minimal SLS realization $\Sigma_m$ of $\Psi$ there exists a surjective SLS morphism
      %$T:\Sigma_r \rightarrow \Sigma_m$. 
      %$\sup_{(\alpha,\beta) n \mbox{selection}, n \in \mathbb{N} } \rank  \hankredu^{H} < +\infty$.
\end{Theorem}

\textcolor{black}{It is well-known that the result above can be made computationally effective when $\mathbb{U}$ and $J$ are finite \cite{MFliessHank,Reut:Book}.}
%SLSs as linear representations
%of recognizable formal power series, and 
%the result above
%and the response maps $H_j$ are identified with their generating series, i.e., with the maps $S_j:Q^{*} \rightarrow \mathbb{Y}$ defined by $S_j(\nu)=H_j(\nu)$ for all $\nu \in Q^{*}$, $j=1,\ldots,d$. In this case, SLSs correspond to linear representations of recognizable formal power series, see e.g.
%  is well-known \cite{MFliessHank,Reut:Book,Son:Real}, and can be made computationally effective.
%In this  case, 
% Theorem \ref{theo:real1}-\ref{theo:real2} are well-known, and can be equivalently strengthened by reducing span-reachability and
%M observability to rank conditions of finite matrices, and 
 %In this case, $W_{\Psi}$ is isomorphic to the column space of
 %the  \emph{Hankel-matrix} $\mathcal{H}_{\Psi}$ defined as follows:
 %for sufficiently large $l,k \in \mathbb{N}$, 
 %which is defined as follows:
 %the condition $\dim W_{\Psi} < +\infty$ by a bounded rank condition
% of finite \emph{Hankel-matrices}. 
 %The Hankel-matrix $H^{\Psi}$ is
 %let    $\nu_1  \prec \nu_2 \prec \cdots$ is an enumeration of $Q^{*}$, and $J=\{j_1,\ldots,j_d\}$ and 
 %\( (\mathcal{H}_{\Psi})_{l,k}=(M(\nu_i\nu_j))_{i,j=1}^{\infty}, \)
 %we can define  Hankel matrix with infinite number of rows and columns as follows:
% 
% \begin{bmatrix} M(\nu_1\nu_1) & \ldots & M(\nu_1\nu_k) \\
% \vdots & \ddots & \vdots \\
% M(\nu_l\nu_1) & \ldots & M(\nu_l\nu_k) \end{bmatrix}
% \] 
 % \in \mathbb{R}^{ld \times kd} \]
 %where $M(v)=$.
 %In particular,  $\dim W_{\Psi}=\Rank~ \mathcal{H}_{\Psi}$, and 
 % can be equivalently replaced by the condition $\sup_{l,k \in \mathbb{N}} \Rank~ H^{\Psi}_{l,k} < +\infty$.
 % and $S_j$ is the generating series of $H_j$, i.e., $S_j(\nu)=H_j(\nu)$ for all $\nu \in Q^{*}$, $j=1,\ldots,d$. 
 %Then $\Psi$ is realizable by a SLS if and only if $\sup_{l,k \in \mathbb{N}} \Rank~ H^{\Psi}_{l,k} < +\infty$, and in this case 
 %Moreover, the canonical realization $\Sigma_{\Psi}$, which is isomorphic to a minimal SLS realization of $\Psi$, 
 %In this case, 
 %a minimal SLS realization can be computed from a finite sub-matrix of $H^{\Psi}$ using a counterpart of the Ho-Kalman algorithm, see e.g. \cite{Son:Real}. 

  %µInspired by this equivalence, in the rest of the paper, 
  %we say that \emph{$H$ has finite Hankel rank}, if $\dim W_{\Psi} <+\infty$, and we call $\dim W_{\Psi}$  the Hankel-rank. 
%A minimal (span-reachable and observable) SLS realization $\Sigma_\Psi$ of $\Psi$ can be constructed from
%$W_\Psi$ as follows \cite[Theorem 2.15]{Son:Real}:
%Let $e_1,\ldots,e_n$ be
%Choose  a basis of $W_\Psi$, and let  $A_u$ be the matrix representation of
%the linear map $W_\Psi \ni \phi \mapsto u \circ \phi$, where $u \circ \phi(v)$ is the shift of $\phi$ by $u$, let  $C$ 
%i.e., $\phi_u(w)=\phi(uw)$ in the basis $e_1,\ldots,e_n$. Let $C$ be 
%be the matrix representation of the linear map $W_{\Psi} \ni \phi \mapsto \phi(\epsilon)$,
%in the basis $e_1,\ldots,e_n$.
%and let $B_j$ be the vector of coordinates of $H_j$ in this basis.
For what follows we need the following technical result.
%MSpan-reachability and observability of SLSs can be characterized as follows.
\begin{Lemma}
  \label{theo:real2}
  %\begin{itemize}
  %\item
  If $\Sigma$ is a minimal SLS realization of $\Psi$, then
  %$\dim \SPAN\{ A_wB_j \mid j \in J, w \in \mathbb{U}^{*}, \|\underline{u}| \le n\}=n$.
  %In particular, 
  there exist sequences $\{v_i,w_i\}_{i=1}^{n} \subseteq \mathbb{U}^{*}$,
  and indices $\{j_i\}_{i=1}^{n} \subseteq  J$
  such that the matrices $O$ and $R$ have rank $n$:
  % $\mathrm{rank}(O)=\mathrm{R}=n$
  \begin{align*}
       R=\begin{bmatrix} A_{v_1}B_{j_1} & \ldots & A_{v_n}B_{j_n} \end{bmatrix} ,\\
      O=\begin{bmatrix} A_{w_1}^TC^T & \ldots & A_{w_n}^TC^T \end{bmatrix}^T\textcolor{black}{.}
  \end{align*}
  %and it is observable 
  %  $\Sigma$ is observable, if and only if 
  %  $\bigcap_{w \in \mathbb{U}^{*}, |w| \le n} \ker CA_w=\{0\}$. In particular,
    %there exists $w_1,\ldots,w_n \in \mathbb{U}^{*}$,
  %and inputs $q_1,\ldots,q_n \in \mathbb{U}$
  %\[ 
  %M\]
  %µ\item 
   %  Any SLS realization of $\Psi$ can be transformed to a span-reachable and observable (minimal)
   %  SLS realization of $\Psi$.
  %\end{itemize}
\end{Lemma}
\begin{proof}
  \begin{color}{black}
  For finite $\mathbb{U}$ the statement is well-known \cite{MFliessHank, Reut:Book,Son:Real}, the proof for the general case is the same, as sketched below. 
 Since $x_{\Sigma}(B_j,w)=A_wB_j$, 
 span-reachability of $\Sigma$ is equivalent to the set $\{A_wB_j \}_{j \in J, w \in \mathbb{U}^{*}}$ containing a basis 
 of $\mathbb{R}^n$.
 By choosing $\{v_i,j_i\}_{i=1}^{n}$ so that $\{A_{v_i}B_{j_i}\}_{i=1}^{n}$ is such a  basis, 
 the first statement follows. As $y_{\Sigma,x_0}(v\sigma)=CA_vx_0$, 
 $v \in \mathbb{U}^{*}, \sigma \in \mathbb{U}$, 
 it follows that
 $y_{\Sigma,x_1}(v\sigma)=y_{\Sigma,x_2}(v\sigma)$ if and only if $x_1-x_2 \in  \ker CA_v$.
 Hence 
 observability of $\Sigma$ is equivalent to $\bigcap_{w \in \mathbb{U}^{*}} \ker CA_w=\{0\}$, i.e., the set of rows of
 $\{CA_w\}_{w \in \mathbb{U}^{*}}$ 
 contains a basis of $\mathbb{R}^{1 \times n}$. By choosing $\{w_i\}_{i=1}^{n}$ so that $\{C_{r_i,\cdot}A_{w_i}\}_{i=1}^{n}$ is such a basis for some output indices $\{r_i\}_{i=1}^{n} \subseteq \{1,\ldots,n_y\}$,
 the second statement follows. 
  \end{color}
% The statements follow by noticing that 
%$\dim \SPAN\{ A_wB_j \}_{j \in J, w \in \mathbb{U}^{*}}=n$
%and observability is equivalent to 
%By choosing a basis 
  %SLSs are equivalent to linear representations of recognizable formal power series.
 %F%or general $\mathbb{U}$ Lemma \ref{theo:real2} is folklore, the closest reference we found is \cite[Section 2]{Son:Real} where a similar result was stated for the so called state-affine systems. The proof of \cite[Section 2]{Son:Real} can easily be
% adapted. % For the sake of completeness its proof is  presented in Appendix.
\end{proof}

 
  %For finite $\mathbb{U}$ and $J$, 
  %the former is  $\sup_{l,k \in \mathbb{N}} \Rank~ H^{\Psi}_{l,k} < +\infty$, and the latter is $\sup_{l,k \in \mathbb{N}} \Rank~ H^{\Psi}_{l,k}$.

% Define a \emph{$n$-selection (selection for short)} as a pair of word-index sets
% $\left(\alpha, \beta \right)$
% such that,
% \begin{equation}
% 	\label{selection:enum}
% 	\begin{split}
% 		&\alpha = \{(\upsilon_{i},k_i)\}_{i=1}^{n}, ~ \beta = \{\nu_l,j_l\}_{l=1}^{n},
% 	\end{split}
% \end{equation}
% for some $\upsilon_{i}, \nu_l \in  Q^{*}$, $|\upsilon_{i}| \le n$, 
% $j_l \in J$, 
% $i,l \in \{1,\ldots,n\}$, $k_i=1,\ldots,p$.
% Define the Hankel matrix $\hankredu^{\Psi} \in \mathbb{R}^{n \times n}$ 
% %and the matrices $\mathcal{H}_{\sigma, \alpha, \beta}^\Psi \in \mathbb{R}^{n \times n}$, $\mathcal{H}_{\alpha}^{\Psi} \in \mathbb{R}^{n \times n_{\mathrm u}}$ and $\mathcal{H}^M_{\beta,q} \in \mathbb{R}^{n_{\mathrm y} \times n}$ as
% follows:
% % in such a way that let us use the ordering of $\alpha$ and $\beta$, and for any 
% %for any $i,j=1,\ldots,n$,
% %the $(i,j)$-th element of %$\hankredu^{M}$ is of the form
% \begin{align}	
% 	& \left[ \hankredu^{H} \right]_{i,l}  \!=\! \left[H_{j_l}\nu_{j_l} \upsilon_{i}) \right]_{k_i}
%    \quad i,l=1,\ldots,n \label{eqn:hankel_u} \\
% \end{align}
% \begin{Theorem}[Recognizable power series \cite{Son:Real}]
%      Assume $\mathbb{U}=Q$ and $J$ are finite.
%       Then $\Psi$ has a realization by a SLS  if and only if 
%       $\sup_{(\alpha,\beta) \mbox{$n$-selection}, n \in \mathbb{N} } \Rank~  \hankredu^{\Psi}=n < +\infty$. Moreover, $\sup_{(\alpha,\beta) \mbox{$n$-selection}, n \in \mathbb{N} } \Rank~  \hankredu^{H}$ is the dimension of any minimal SLS realization of $\Psi$. 
% \end{Theorem}
% In this case  $\dim W_{\Psi}=\sup_{(\alpha,\beta) \mbox{$n$-selection}, n \in \mathbb{N}} \Rank~  \hankredu^{\Psi}$, and  if $\dim W_{\Psi} < +\infty$, then $W_{\Psi}$ is isomorphic to $\mathrm{Im} \hankredu^{\Psi}$ for a suitable selection $(\alpha,\beta)$.
% %Note that $\dim W_H=\sup_{(\alpha,\beta) \mbox{$n$-selection}, n \in \mathbb{N} } \Rank~  \hankredu^{H}$ and in fact a minimal SAS realization can be computed as follows. 
% Moreover, we can compute a minimal SLS of $\Psi$ from 
% $\hankredu^{\Psi}$ as follows. 
% %\cite{Son:Real}, the generating series $S_H$ is uniquely determined by $H$.
% %The algorithm is as follows. 
% Let us define the shifted Hankel matrices
% \begin{align}
% & \left[\mathcal{H}_{q, \alpha, \beta}^{\Psi}\right]_{i,l} = \left[ H_{j_l}(\nu_{l}q\upsilon_{i}) \right]_{k_i} , ~i,l=1,\ldots,n \label{eqn:hankelu:shift} \\
% 	& \left[\mathcal{H}^{\Psi}_{\alpha,j}\right]_{i} \! = \! \left[ H_j(\upsilon_{i})\right]_{k_i},  \ i\!=\!1,\ldots,n \label{eqn:hankel_u_2a}\\
% 	& \left[\mathcal{H}^{\Psi}_{\beta}\right]_{l} \!=\! \left[ H_{j_l}(\nu_{l}) \right],  l=1,\ldots n
%  \label{eqn:hankel_u_2b}
% \end{align}
% \begin{algorithm}[h]
% 	\caption{Ho-Kalman for SAS}
% 	\label{algo:Nice_select_true_deter_basic}
% 	\textbf{Input}:  $n$-selection $\left(\alpha, \beta \right)$,  and $\Psi$
% 	\hrule\vspace*{.1cm}
%       Construct the matrices 
%       $\mathcal{H}^\Psi_{\alpha,\beta}$, $\mathcal{H}^\Psi_{q, \alpha, \beta} $ %\eqref{eqn:hankelu:shift}; Hankel matrices   
%      $\mathcal{H}_{\alpha,\j}^\Psi$ and $\mathcal{H}_{\beta}^\Psi$, and set
% 		\begin{align*}
% 			&  \hat{A}_{q} = (\hankredu^\Psi)^{-1}  \mathcal{\Psi}_{q, \alpha, \beta}^H, ~  \hat{B}_j = (\hankredu^\Psi)^{-1} \mathcal{H}^\Psi_{\alpha,j}, ~ \\
%       & \hat{C} =  \mathcal{H}^\Psi_{\beta}
% 		\end{align*}
% 	\hrule 
% 	\textbf{Output}: SLS $(n,\{\hat{A}_{q}\}_{q  \in Q},\hat{B},\hat{C})$
% \end{algorithm}
% \vspace{-0.1cm}
% \begin{Lemma}[Ho-Kalman \cite{Son:Real}]
% 	\label{basis_red:lemma}
% 	%Let , and 
% Assume that $n=\sup_{(\alpha,\beta) \mbox{$n$-selection}, n \ge 1 } \Rank~ \hankredu^{\Psi}$.
% Then there exists an $n$-selection $(\alpha,\beta)$ such that $\text{rank}(\hankredu^\Psi) = n$.
%         Algorithm \ref{algo:Nice_select_true_deter_basic}
%         returns a minimal SLS realization of $\Psi$, which is imorphic to the canonical
%         realiztaion $\Sigma_\Psi$. 
% \end{Lemma}



     %We call $\beta$ \emph{decay rate}.
%%     The family $\Psi$ is \emph{uniformly absolutely summable (UAS)}, 
%%     %if there exists $K > 0$ such that 
%%     %for every infinite sequence $w \in \mathbb{U}^{\infty}$, $j \in J$
%%      if the series
%%     \( \sum_{s=1}^{\infty} \|H(w[i] \cdots w[1])\|\) 
%%     are uniformly bounded and convergent over $j \in J$, $w \in Q^{\infty}$, i.e.
%%     \[
%%       \begin{split}
%%       &  \sup_{j \in J, w \in Q^{\infty}} \sum_{s=1}^{\infty} \|H(w[i]\cdots w[1])\|_2 < +\infty \\
%%        & \lim_{N \rightarrow \infty}\sup_{j \in J, w \in Q^{\infty}} \sum_{s=N}^{\infty} \|H(w[i]\cdots w[1])\|_2=0
%%       \end{split}
%%     \]
     %for every $\epsilon > 0$ there exists $N_{\epsilon}$ such that for all
     %$N > N_{\epsilon}$, for all 
     %$j \in J$, $w \in \mathbb{U}^{\infty}$, 
     %$\sum_{s=N}^{\infty} \|H(w[i]\cdots w[1])\|_2 < \epsilon$.
     %$w \in \mathbb{U}^{\infty}$, $j \in J$ the serie
     %
     %s
     %\( \sum_{i=1}^{\infty} \|H(w[1] \cdots w[i])\| \) are uniformly convergent, % and the 
%     convergence is uniform, 
     %i.e.
     %for every $\epsilon > 0$ there exists $N_{\epsilon}$ such that for all
     %$N > N_{\epsilon}$, for all 
     %$j \in J$, $w \in \mathbb{U}^{\infty}$, 
     

%Clearly, if $\Sigma$ is UGES with decay rate $\beta$ if and only if 
%any solution of \eqref{eq2} it holds that $\|x(t)\| < C\beta^t \|x(0)\|$ and it its
%uniformly globally exponentially stable dynamical system \cite{Son:MathContr}.

 %does not depend on $u(t)$, $H(u(0)\cdots u(t))$ depends only on $u(0),\ldots, u(t-1)$, if
 %$H$ has a realization by SAS; if we include  a term in the readout map which depends on $u(t)$,
 %i.e. we set $y(t)=H(u(t))x(t)+d(u(t))$, then $H(u(0)\cdots u(t))$ will depend on $u(t)$. 
 %The results of this paper can readily be extended to that more general case at the 
 %expense of more cumbersome notation.
 %\end{Remark}
%$\sum_{i=N}^{\infty} \|H(w[1]\cdots w[i])\|_2 < \epsilon$.
    % \subsection{Stability and exponential decay}

%     We say that $\Psi$ is UED, UAS, UD respectively, if each element of $\Psi$
%     is UED, UAS, UD respectively. 
%\begin{Remark}[Intuition for exponential forgetting]
%    Exponential forgetting captures the idea that the effect of past inputs 
%    decays exponentially over time. Specifically, if we apply the same input sequence 
%    $u_1 \cdots u_k$ but with different "prefixes" $w$, the difference in outputs 
%    vanishes exponentially as $k$ increases. This property is fundamental for stable systems, 
%    as it ensures that the system's output depends primarily on recent inputs rather than 
%    the entire input history. For state-affine systems with contracting dynamics, 
%    exponential forgetting follows naturally from the decay of the state trajectory.
%\end{Remark}



\section{Main result}
We call the SLS  \eqref{eq1} \emph{globally uniformly  exponentially stable (GUES)
with decay rate $\beta$}, if \textcolor{black}{the set of initial states $B=\{B_j\}_{j \in J}$ is bounded} and
  there exists $C_g \ge 1$ such that 
    \begin{equation}
    \label{def:sls:gues}
    \textcolor{black}{\forall w \in \mathbb{U}^{*}: }\quad 
    \|A_w\|_2 \le C_g\beta^{|w|}. 
    \end{equation}
Note \textcolor{black}{that} \eqref{def:sls:gues} \textcolor{black}{implies that} any solution of \eqref{eq1} \textcolor{red}{satisfies} $\|x(t)\|_2 < C\beta^t \|x(0)\|_2$ \cite{RJungers}. 
%By \cite{RJungers}, GUES is equivalent to
%the joint spectral radius of $\{A_u\}_{u \in \mathbb{$
%$\Sigma$ is GUES with decay rate $\beta$ if and only if 
    The family $\Psi$ from \eqref{family:resp} 
     is called \emph{uniformly decaying} (UD), if 
     $\sup_{\substack{w \in \textcolor{black}{\mathbb{U}^{*}}, \\ j \in J}} \|H_j(w)\|_{\textcolor{black}{2}} < +\infty$, and 
     \begin{equation}
     \label{def:ud}
     \lim_{t \rightarrow \infty} \sup_{w \in \mathbb{U}^{*}, |w|=t,  j \in J} \|H_j(w)\|_{\textcolor{black}{2}}=0 
     \end{equation}
     and $\Psi$ is called  \emph{uniformly exponentially decaying (UED) with decay rate $\beta \in (0,1)$}, 
     if there exists $C_{ed} \ge 1$, %$\beta \in (0,1)$
     such that %$v \in \mathbb{U}^{*}$ 
     \begin{equation}
     \label{def:ued}
     \forall w \in \mathbb{U}^*: ~ \sup_{j \in J} \|H_j(w)\|_{\textcolor{black}{2}} \le C_{ed}\beta^{|w|}.
     \end{equation}
     Clearly, UED implies UD.
     Now, we relate realizability by a stable SLS to these properties.
     %the main result of the paper, which states that the converse also holds
     %for families realizable by SLSs. 
  \begin{Theorem}
  \label{main:theo}
  For a family $\Psi$ as in \eqref{family:resp}, the following hold.
  \\
    \textbf{(A)}
    $\Psi$ can be realized by a SLS which is GUES 
    if and only if $\Psi$ is UD  and $\Psi$ has a finite Hankel-rank.
    \\
    %\item{\textbf{(B)}}
    %The family of response map $\Psi$ can be realized by a SLS which is GUES with decay rate $\beta$ 
    %if and only if $\Psi$ is UED with rate $\beta$ and 
    %$\Psi$ has a finite Hankel-rank.
    \textbf{(B)}
    If $\Psi$ is UED with rate $\beta$, then all minimal SLS realizations of $\Psi$ are GUES with rate $\beta$.  
    \\
  \textbf{(C)}
      If $\Psi$ has a finite Hankel-rank and $\Psi$ is UD, then $\Psi$ is also UED.     
   % \end{itemize}
  \end{Theorem}
  \begin{proof}%[Proof of Theorem \ref{main:theo}]
     %Throughout the proof, for $w \in \mathbb{U}^{*}$, $F(w)$ stands for the matrix product $A_w$ as
     %in Notation \ref{mat:prod} with $A_u=F(u)$, $u \in \mathbb{U}$.
     Let $\Sigma$ \eqref{eq1}
     be a realization of $\Psi$.
     %Assume that $\Sigma$ is a SLS of the form \eqref{eq1} and it is a realization of $\Psi$. Then
     %\[
     %  \begin{split}
     %   & H(u_0\cdots u_t)= d_{u(t)} + C_{u(t)}b_{u(t-1)}+ \\
     %   & \sum_{s=0}^{t-2} C_{u(t)}F_{u(t-1)\cdots u(s+1)}b_{u(s)} 
     %  \end{split}
     %\]
     %and hence 
     %for every $w \in \mathbb{U}^{*}, j \in J$ %u,r \in \mathbb{U}$
     %\(
     %    H_j(w)=CA_w B_j 
     %M\).
     The theorem follows from the subsequent implications.
     
    \textbf{(I) $\Sigma$ is GUES with rate $\beta$ $\implies$ \eqref{def:ued} holds.}
      %Assume that $\Psi$ can be realized by a SAS which is GUES with decay rate $\beta$.
    % the joint spectral radius of $\{F(u) \mid u \in \mathbb{U}\}$ is smaller than $\beta$,
    % \cite{JungersBook}.
     %Since %$\Sigma$ is GUES, for some $K > 0$, 
     %$\|A_w\|_2 < K \beta^{|w|}$. %From \eqref{eq:real2} it follows that
     %since \( H_j(w)=CA_wB_j\) and 
     %for all $w=v_1\cdots v_t \in \mathbb{U}^{+}$, $t \ge 1$,
     %\[ 
     %  \begin{split}
     %    & H(qvu_1 \cdots u_k)-H(u_1\cdots u_k)  = \\
     %   & \sum_{i=1}^{t-1} H_{u_k}F_{u_1\cdots u_{k-1}}F_{v_{i+1}\cdots v_t}b_{v_i}
     %  \end{split}
     %\]
    %and as 
    %\[ \|CA_wB_j\|_2 \le K \|C\|_2 \sup_{j \in J}  \|B_j\|_2 \beta^{|w|} \]
   % If \textbf{(A)} holds, then $\sup_{u \in \mathbb{U}} \|b(u)\|_2$ is well defined
   % as $\b(u)$ is a polynomial function of $u$ and $\mathbb{U}$ is bounded.
   % If \textbf{(B)} holds, then $\|b(u)\|_2 \le \sum_{i=1}^{n_p} \psi_{i}(u) \|b_i \|_2$ and 
   % by assumption $\psi_i(u) \in [0,1]$ and $\sum_{i=1}^{n_p} \psi_i(u) \le 1$ and hence
   % $\sup_{u \in \mathbb{U}} \|b(u)\|_2=\le \max_{q \in Q} \|b_q\|$. 
   Since 
   $\|H_j(w)\|_2=\|CA_wB_j\|_2
   \le \|C\|_2 \|A_w\|_2\|B_j\|_2$,
   if
   \eqref{def:sls:gues} holds, then
    with $C_{ed}=C_g \|C\|_2 \sup_{j \in J} \|B_j\|_2$ \eqref{def:ued}
    holds.
    %\(
     %  \|H_j(w)\|_2=\|CA_wB_j\|_2 \le K_c\beta^{|w|}
      %\begin{split}
      % & \|H(qvu_1 \cdots u_k)-H(u_1\cdots u_k)\|_2 \le \\
      % & \sum_{i=1}^{t-1} \|C_{u_k}F(u_1\cdots u_k)F(v_{i+1}\cdots v_t)b(v_i)\|_2 \le  \\
      % & \sum_{i=1}^{t-1} \sup_{u \in \mathbb{U}} \max\{\|C_u\|,\|b(u)\|_2\} \beta^{i+k-1} \le
      % C \beta^{k} 
      %\end{split}
    %\) holds for all $w \in \mathbb{U}^{*}$.
    %That is $\Psi$ is UED.
    % and hence it is UAS and UD.
    
  \textbf{\textcolor{red}{(II)} $\Psi$ is UD $\implies$ all minimal SLSs are GUES.}
    %exponentially decaying with decay rate $\beta$ and 
    %$\sup_{(\alpha,\beta) \mbox{$n$-selection}, n \ge 1 } \Rank~ \hankred^. Then by Theorem \ref{theo:real1}-\ref{theo:real2}, there exists
    \textcolor{red}{Let} $\Sigma$ be  a minimal SLS which realizes $\Psi$.
    Then the matrix $O$ from Lemma \ref{theo:real2} has a left inverse $O^{+}$, and 
    the matrix $R$ from Lemma \ref{theo:real2} is invertible.
    %there exists sequences $w_1,\ldots, w_n \in \mathbb{U}^{*}$ such that the matrix
    %\[  O = \begin{bmatrix} (CA_{w_1})^T, (CA_{w_2})^T,\ldots,(CA_{w_n})^T \end{bmatrix}^T
    %\]
    % is full column rank, i.e.,  left inverse $O^{+}$ exists, and 
    %\textbf{Case of (A)} By \cite[page 147]{SonPoly76} there exists 
    % sequences $w_1^{*},\ldots, w_t^{*} \in (\mathbb{R}^{n_p})^{*}$ such that
    %\[  O(w_1^{*},\ldots, w_t^{*}) = \begin{bmatrix} (HF(w_1^{*}))^T, (HF(w_2))^T,\ldots,(HF(w_t^{*})^T \end{bmatrix}^T
    %\]
    %has zero kernel. 
    %Here, as $F:\mathbb{U} \ni  u \mapsto F(u)$ is a  polynomial function, 
    % we can extend it from 
    %$\mathbb{U}$ to $\mathcal{R}^{m}$ uniquely, and 
    %for any $w \in (\mathbb{R}^{m})^*$, $F(w)=A_w$ as in Notation \ref{matrix:prod} for $A_u=F(u)$. 
    %Let the length of $w_i^{*}$ to be $t_i$. Moreover, since $O(w_1^{*},\ldots, w_t^{*})$ has $n$
    %columns, we can choose $n$ elements of $w_1^{*},\ldots, w_t^{*}$ such that the matrix
    %formed by the block rows $O(w_1^{*},\ldots, w_t^{*})$ corresponding to those 
    %elements is of rank $n$, i.e., we can assume that $t \le n$. 
    %Every sequence $w_i (\mathbb{R}^{m})^{*}$ of length $t_i$ can be viewed as an element
    %of $()\mathbb{R}^m)^{t_i}$ and we can define the map
    %\[
    %  \begin{split}
    %   O:(\mathbb{R}^m)^{t_1} \cdots (\mathbb{R}^m)^{t_n} \ni (w_1,\ldots, w_n) 
    %   \mapsto \\
    %   \begin{bmatrix} (HF(w_1))^T, (HF(w_2))^T,\ldots,(HF(w_n)^T \end{bmatrix}^T
    %  \end{split}
    %\]
    %Since the entries of $O(w_1,\ldots, w_n)$ are polynomial in $w_1,\ldots,w_n$,
    %and there exists a $n \times n$ submatrix of $O(w_1^{*},\ldots,w_n^{*})$ whose
    %determinant is non-zero, it means that for almost all $w_1,\ldots,w_n$, 
    %the determinant of the corresponding $n \times n$ sub-matrix of $O(w_1,ldots,w_n)$ is not
    %zero, i.e., for almost all $w_1,\ldots,w_n$, $O(w_1,ldots,w_n)$ has zero kernel.
    %But as $\mathbb{U}$ has non-empty interior, it cannot be a set of measure zero, and 
    %hence $(\mathbb{U})^{t_1} \times \cdots \times (\mathbb{U})^{t_n}$ is not a set of measure
    %zero in $(\mathbb{R}^m)^{t_1} \cdots (\mathbb{R}^m)^{t_n}$. That is, $w_1^{*},\ldots, w_n^{*}$
    %can be chosen from $\mathbb{U}^{*}$. 
%    Consider now the dynamica system
%    \[ x(t+1)=F(u(t))x(t), ~ x(0)=x_0 \]
    %From span-reachability it follows that for some
    %there exists  $v_1,\ldots, v_k \in \mathbb{U}^{*}$, 
    %$j_1,\ldots, j_k \in J$, 
    %such that the matrix 
    %\[ R=\begin{bmatrix} A_{v_1} B_{j_1}& \ldots  & B_{v_n}B_{j_n} \end{bmatrix} \]
    %has rank $n$, i.e., it 
    %is invertible. 
    %$x_0=\sum_{i=1}^{k} \alpha_i F(v_i)b_{q_i}$. Then 
    %\[
    %   x(t)=\sum_{i=1}^k F(\underline{u})F(v_i)b_{q_i} 
    %\]
      %O(w_1,\ldots,w_n)x(t)= \sum_{i=1}^{k}
      % \alpha_i \begin{bmatrix} HF(w_1^{*})F(\underline{u})F(v_i)b_i \\
      %    HF(w_2^{*}F(\underline{u})F(v_i)b(q_i) \\
      %   \vdots,
      %   HF(w_n)F(\underline{u})F(v_i)b(q_i) \end{bmatrix}
    % \]
    \textcolor{red}{For} any $\underline{u} \in \mathbb{U}^{*}$,
    %\begin{equation}
    %\label{main:theo:pf1} 
     \( CA_{w_l}A_{\underline{u}}A_{v_i}B_{j_i} \!\!=\!\! H_{j_i}(v_i\underline{u}w_l) \), %H_{j_i}(v_i\underline{u}w_l) \), and
    %\end{equation}
    %Hence, 
    %\begin{equation}
    %\label{main:theo:pf2} 
    %  \begin{split}
    %\|CA_{w_l}A_{\underline{u}}A_{v_i}B_{j_i}\|_2  = \|H_{j_i}(v_i\underline{u}w_l)\|_2
    %  \end{split}
    %\end{equation}
    hence 
    \begin{equation*}
    %\label{main:theo:pf3} 
      \begin{split}
      & \|OA_{\underline{u}}R\|_2 
      %\|OA_{\underline{u}}R\|_F \le \\
      \le \sum_{i,l=1}^{n} \|CA_{w_l}A_{\underline{u}}A_{v_i}B_{j_i}\|_2\!\!= \!\!
         \sum_{i,l=1}^{n}  \|H_{j_i}(v_i\underline{u}w_l)\|_2 %H_{j_i}(v_i\underline{u}w_l)\|_2
      \end{split}  
    \end{equation*}
    %for all $\underline{u} \in \mathbb{U}^{*}$.
    %\textcolor{pink}{FLORIAN: I removed an equality here because it was missing a matrix $R$ but now then became too long to fit nicely into the column (see the LaTeX).}\\
    As
    %$A_{\underline{u}}=O^{+}OA_{\underline{u}} \textcolor{black}{R} R^{-1}$ and 
    $\|A_{\underline{u}}\|_2\!\!=\!\!\|O^{+}OA_{\underline{u}} \textcolor{black}{R} R^{-1}\|_2 \le  \|O^{+}\|_2 \|OA_{\underline{u}}R\|_2 \|R^{-1}\|_2$,
    \begin{equation}
    \label{main:theo:pf4}
    \begin{split}
     & \|A_{\underline{u}}\|_2
     %=\|O^{+}OA_{\underline{u}}RR^{-1}\|_2 
     %\le  \|O^{+}\|_2 \|OA_{\underline{u}}R\|_2 \|R^{-1}\|_2 
     \le \|O^{+}\|_2 \|R^{-1}\|_2 
        \sum_{i,l=1}^{n}  \|H_{j_i}(v_i\underline{u}w_l)\|_2.
    \end{split}
  \end{equation}
  If $\Psi$ is UD, then by setting $S=\sup_{w \in \mathbb{U}^{*}, j \in J} \|H_j(w)\|_2$
  it follows from \eqref{main:theo:pf4} that 
  %\begin{equation} 
  %    \label{main:theo:pf5}
  \textcolor{black}{$\|A_{\underline{u}}\|_2 \le \|O^{+}\|_2 \|R^{-1}\|_2 n^2S$} 
  %\sum_{i,l=1}^{n}  \|H_{j_i}(v_i\underline{u}w_l)\|_2 \}$
  %\end{equation}
  i.e., the set $\{ A_{u} \}_{u \in \mathbb{U}}$ is bounded. 
   \textcolor{black}{From} \eqref{def:ud}
   %$\|H_{j_i}(v_i\underline{u}w_l)\|_2$ converges to zero 
   %uniformly as $|\underline{u}| \rightarrow +\infty$, 
   it follows that
    for every $\epsilon > 0$ there exists $N_{\epsilon}$ such that
    if $|\underline{u}| > N_{\epsilon}$, \textcolor{black}{then} 
    \[  \begin{split}
    %  & \|CA_{w_l}A_{\underline{u}}A_{v_i}B_{j_i}\|_2  = \\
       \|H_{j_i}(v_i\underline{u}w_l)\|_2 < \frac{\epsilon}{n^2 \|O^{+}\|_2 \|R^{-1}\|_2 }
      \end{split}.
    \]
  Hence, by using \eqref{main:theo:pf4} it follows that 
  \( \|A_{\underline{u}}\|_2 \le \epsilon \)
  \textcolor{black}{if} $|\underline{u}| > N_{\epsilon}$, i.e.,
  \eqref{eq1} is globally uniformly asymptotically stable. 
  It is well-known that then \eqref{eq1} is \textcolor{black}{GUES}
  % globally uniformly exponentially stable
  \cite[Theorem 2.15]{Sun:Book}, \cite{Piatnitsky}.
  %\footnote{It is well-known, e.g., \cite[Theorem 2.15]{Sun:Book}, \cite{Piatnitsky} that under mild assumptions (for instance $\mathbb{U}$ being finite) is  uniform asymptotic stability of \eqref{eq2} implies uniform global exponential stability.  However, we have not found in the literature a formal statement which precisely covers, although the available proofs seem to extend to our case. For this reason, we preferred to complete the argument.}. 
  %In particular, it the follows that the joint spectral radius of 
  %Hence, any solution $x$ of \eqref{eq2} satisfies 
  %$\lim_{t \rightarrow \infty} x(t)= 0$. 
  %For the sake of completeness, the proof is as follows: as 
  Indeed, as $\{ A_u\}_{u \in \mathbb{U}}$ is bounded, by
  \cite[Corollary 1.1]{RJungers} 
  the joint spectral radius
  $\rho=\lim_{t \rightarrow \infty}\sup_{\underline{u} \in \mathbb{U}^t} \|A_{\underline{u}}\|_2^{1/t}$ of   \textcolor{black}{$\{A_u\}_{u \in \mathbb{U}}$ satisfies $\rho \!\!<\!\! 1$.}
  Let $\epsilon > 0$ be such that $\beta=(\rho+\epsilon)<1$.
  Then there exists $T_{\epsilon}$ such that if $|\underline{u}| > T_{\epsilon}$
  then $\|A_{\underline{u}}\|_2 < \beta^{|\underline{u}|}$.
 % By setting 
  \textcolor{black}{For \textcolor{black}{$|\underline{u}| \le T_{\epsilon}$},
  $\|A_{\underline{u}}\|_2 \le \textcolor{black}{n^2S\|O^{+}\|_2 \|R^{-1}\|_2} \textcolor{black}{\le n^2S\|O^{+}\|_2 \|R^{-1}\|_2 \beta^{|\underline{u}|-T_{\epsilon}}}$.
  Hence, 
  \eqref{def:sls:gues} holds with 
  $C_g\!\!=\!\!\max\{1, \textcolor{black}{n^2S\|O^{+}\|_2 \|R^{-1}\|_2 \beta^{-T_{\epsilon}}}\}$.}
  %\( \|A_{\underline{u}}\|_2 \le K\beta^{|\underline{u}|}. \)
  %As joint spectral radius is closed under taking closure 
  %bounded families of matrices \cite[Proposition 1.8]{JungersBook}, 
  %by \cite[Corollary 1.1]{JungersBook}
  %then 
  %using either \cite[Theorem 2.15]{SunBookStab} or 
  %By using boundedness of $\{A_u \mid u \in \mathbb{U}\}$
  %from \cite{Pyatnitski} (or by repeating the proof of \cite[Theorem 2.15]{SunBookStable}) 
  %we get that  
  %\eqref{eq2} is uniformly exponentially stable. The latter
  %(by taking the difference inclusion generated by the convex combination of $\{A_u\}_{u \in \mathbb{U}}$)
  %implies that there exists $C > 0$, $\beta \in (0,1)$ such that 
   %and for $\underline{u} \in \mathbb{U}^{*}$ such that $|\underline{u}|  > N_{\epsilon}$, 
   %   \[ 
   %   \begin{split}
   %   & \|OA_{\underline{u}}R\|_2 =
   %   \|OA_{\underline{u}}R\|_F \le \\
   %   & \sum_{i,l=1}^{n} \|CA_{w_l}A_{\underline{u}}A_{v_i}B_{j_i}\|_2 \le \epsilon
      %\le K_1 \beta^{t}
    %  \end{split}  
    %\]
    \textcolor{black}{Also},
   \(
    %\begin{split}
   %& 
   \|OB_j\|_2 \le 
   %& \|\begin{bmatrix} (CF_{w_1}B_j)^T, (CF_{w_2}B_j)^T,\ldots,(CF_{w_n}B_j)^T \end{bmatrix}^T\|_2 \le \\
     \sum_{i=1}^{n} \|CA_{w_i}B_j\|_2=
      \sum_{i=1}^{n} \|H_j(w_i)\|_2 \le \textcolor{black}{nS}
      %& \le \sum_{i=1}^{n} C\beta^{|w_i|}=:K_b
   %\end{split}
   \)
   and hence
   \( \|B_j\|_2=\|O^{+}O B_j\|_2 \le \|O^{+}\|_2 \textcolor{black}{nS}  \).  % K_b
    %Moreover, for any $u \in \mathbb{U}$
    %\[ 
    %  \begin{split}
    %   & C_uR=\begin{bmatrix} C_uF_{v_1} b_{q_1}& \ldots  &C_uF_{v_n}b_{q_n} \end{bmatrix} 
    %   \\   
    %   & =\begin{bmatrix}  H(q_1v_1u)-H(v_1u) & \ldots & H(q_nv_nu)-H(v_nu) \end{bmatrix}
    %  \end{split}
    %\]
    %and hence 
    %\[
    %  \begin{split}
    %   & \|C_uR\|_2 \le \|C_uR\|_F \le \sum_{i=1}^{k} \|H(q_1v_1u)-H(v_1u)\| \\
    %   & \le \sum_{i=1}^{n} C \beta^{2+|v_i|}=:K_C
    %  \end{split}
    %\]
    %and hence
    %\[ \|C_u\|_2=\|CRR^{-1}\|_2 \le \|C_u R \|_2\|R^{-1}\| \le K_C \|R^{-1}\|_2 \]
   %Similarly, 
   
\textbf{(III) If $\Sigma$ is minimal and \eqref{def:ued} holds, then \eqref{def:sls:gues} holds with the same $\beta$}. 
  It follows from \eqref{def:ued} that
  \( \|H_{j_i}(v_i\underline{u}w_l)\|_2 <  C_{ed}\beta^{|v_i|+|\underline{u}|+|w_l|} \).
  Using \eqref{main:theo:pf4}, 
  %\( \|A_{\underline{u}}\|_2 \le K\beta^{|\underline{u}|} \)
  %where 
  and choosing \textcolor{black}{$C_g=\|O^{+}\|_2\|R^{-1}\|_2\max\{1,C_{ed}\sum_{l,i=1}^{n} \beta^{|v_i|+|w_l|}\}$},  it follows that
  \eqref{def:sls:gues} holds. 
  %Since UED implies UD, the other properties of GUES SLS follow in the same
  %way as above.  
%%  \[ 
%%    \begin{split}
%%     \|O(w_1,\ldots,w_n)x(t)\|_2 \le
%%     \| \sum_{i=1}^{k}
%%       \alpha_i \begin{bmatrix} HF(w_1)F(\underline{u})F(v_i)b_i \\
%%          HF(w_2F(\underline{u})F(v_i)b(q_i) \\
%%         \vdots,
%%         HF(w_n)F(\underline{u})F(v_i)b(q_i) \end{bmatrix} \|_2 \leq \\
%%         \sum_{i=1}^{k}
%%       |\alpha_i| \|\begin{bmatrix} HF(w_1)F(\underline{u})F(v_i)b_i \\
%%          HF(w_2F(\underline{u})F(v_i)b(q_i) \\
%%         \vdots,
%%         HF(w_n)F(\underline{u})F(v_i)b(q_i) \end{bmatrix} \|_2 \le
%%      \sum_{i=1}^{k}
%%      \sqrt{n} \sum_{j=1}^{n} 
%%    \|HF(w_j)F(\underline{u})F(v_i)b(q_i)\|
%%      \le \underbrace{(C \sum_{i=1}^k \sum_{j=1}^{n} \beta^{|v_i|+|w_j|})}_{=:K}\beta^{t}
%%      =K \beta^{t}
%%    \end{split}
%%  \]
%%  That is, as $x(t)=O^{+}(w_1,\ldots,w_n)O(w_1,\ldots, w_n)x(t)$ and 
%%  \[ \|x(t)\| \le \|O^{+}(w_1,\ldots,w_n)\|_2\|\|O(w_1,\ldots,w_n)x(t)\|_2 \le 
%%      \tilde{K}\beta^{t}  
%%  \]
%%  with $\widetilde{K}=K\|O^{+}(w_1,\ldots,w_n)\|_2$.
%%  From this it follows that $\lim_{t \rightarrow \infty} x(t)=0$. Then by 
%%  \cite[Corollary 1.1]{JungersBook} 
%%  $\{F(u) \mid u \in \mathbb{U}\}$ has a joint spectral radius which is smaller than $1$,
%%  which implies that 

  %Finally, \textbf{(D)} part of the theorem follows by recalling from Theorem \ref{theo:real1} that 
  %if $\dim W_{\Psi} <+\infty$ if and  only if $\Psi$ has a SLS realization, and that if $\Psi$ has a SAS realization then it has a minimal realization.
  %Since in
  %the proof we show that $\Psi$ being UD implies that \eqref{eq2} is uniformly asymptotically
  %stable, then we could reuse 
  %\cite[Theorem 2.15]{SunBookStab}, \cite{Pyatnitski} to simplify the proof whenever 
  %the corresponding mild assumptions are satisfied. This is the case for instance one
  %mathbb{U}$ is finite.  However, it is
  Statement \textbf{(A)} follows from \textbf{(I)}, \textbf{(II)}, and the fact that finite Hankel-rank guarantees the existence of a minimal SLS realization. Statement \textbf{(B)} follows from \textbf{(III)}.  Statement \textbf{(C)} follows from \textbf{(II)} and \textbf{(I)}.
  \end{proof}


\section{Special cases} %: state-affine systems}

In this section we apply the main results to  SASs, LPV-SSAs and LSSs, by associating any particular system of those classes with a suitable SLS following \cite{Son:Real},\cite{Petreczky2016}.
%In comparison to SLS for these systems a natural assumption is that the initial state is zero. 

\subsection{State-affine systems}
  %Consider a state-affine system  
 %µ\begin{equation}
  %  \label{eq1:sas} 
  %     \begin{split}
  %       &  x(t+1)=F_{u(t)}x(t) + b_{u(t)}, \quad x(0)=0\\
  %       & y(t)=Cx(t)
       %\end{split}
   %\end{equation}
   %where $F_{u(t)} \in \mathbb{R}^{n \times n}$, $b_{u(t)} \in \mathbb{R}^{n}$, 
   %$C_{u(t)} \in \mathbb{R}^{p \times n}$, $d_{u(t)}$, 
   % $x(t) \in \mathbb{R}^n$ is the state, $u(t) \in \mathbb{U} \subseteq \mathbb{R}^{m}$ is
   %the input and $y(t) \in \mathcal{Y} \subseteq \mathbb{R}^p$.
  We consider SASs of the form \eqref{eq1:sas}. These systems are SO, and their %corresponding
  \textcolor{black}{response maps are of the form} $H:\mathbb{U}^{*} \rightarrow \mathbb{R}^{\textcolor{black}{n_y}}$.
  Following \cite{Son:Real}, we can reduce realization theory of SASs to realization theory of SLSs as follows. 
  %Assume that $H$ is an input-output map which is known to
  % be a response map and for every $u \in \mathbb{U}^{+}$
   We associate with $H$ the 
   family $\Psi_H=\Psi$ of the form
   \eqref{family:resp}, where 
   $J=\mathbb{U}$ and 
   \begin{equation}
   \label{def:hu}
     \forall w \in \mathbb{U}^{*}, u \in \mathbb{U}: \quad H_u(w)=H(uw)-H(w) \textcolor{black}{.}
   \end{equation}
   %Let us response map $H_u:\mathbb{U}^{*} \rightarrow \mathbb{R}^p$ such that
    %$H_u(w)=H(uw)-H(w)$ for $w \in \mathbb{U}^{+}$, and $H_u(\epsilon)=H(u)$. 
    Clearly, 
    $H(w)\!=\!\textcolor{black}{H(\epsilon)+\sum_{s=1}^{|w|}} H_{w[s]}(w[s+1]\cdots w[|w|])$, \textcolor{black}{and}
    \eqref{eq1:sas} is a realization of $H$, if and only if for all 
    $v \in \mathbb{U}^{*}$
    \[ 
    H_u(v)=CA_v (A_ux_0+b_u-x_0),  ~ u \in \mathbb{U}, ~ H(\epsilon)=Cx_0+d. 
    \]
    %and $d=H(\epsilon)$. 
    % \begin{equation} 
    % \label{sas:io:eq1}
    %    H(w)= \sum_{s=0}^{t-1} CA_{w[s+1]\cdots w[t-1]}b_{w[s]}=\sum_{s=0}^{t-1} H_{w[s]}(w[s+1] \cdots w[t-1])  
    %µ \end{equation}
    %and 
     %Hence, by \cite{Son:Real} 
     %$H$ 
     %that a input-output map $H:\mathbb{U} \rightarrow \mathbb{R}^p$ has a 
     %realization by a SAS \eqref{eq1:sas}, if and only if 
     %the family $\Psi_H=\{ H_u \mid u \in \mathbb{U}\}$ of response maps with the index set 
     %$J=\mathbb{U}$ has a realization by a SLS. Moreover, 
    Following \cite{Son:Real}, if
    $\Sigma$ is as in \eqref{eq1:sas}, we  define
    the SLS
    \[
    R_{\Sigma}=\bigl(\textcolor{black}{\{A_u\}_{u\in\mathbb{U}}},B=\{B_u:=A_ux_0+b_u-x_0\}_{u\in\mathbb{U}},C)\textcolor{black}{.}
    \]
    %whose index set is $J=\mathbb{U}$.  %For $H$ we define the family of response maps
    %$\Psi_H=\{H_u\midi}_{u \in \mathbb{U}\}$. 
    Then $\Sigma$ is a realization of $H$
    if and only if $R_{\Sigma}$ is a SLS realization of $\Psi_H$
    and $Cx_0+d=H(\epsilon)$.
    \textcolor{black}{The} transformation
    $\Sigma \mapsto R_{\Sigma}$ preserves observability and isomorphism, and maps affine-reachable SASs to span-reachable SLSs. 
    Conversely, if
    $R=(\{A_{u}\}_{u \in \mathbb{U}},B,C)$
    is a SLS realization of
    $\Psi_H$, then 
    the SAS $\Sigma_R$ of the form
    \eqref{eq1:sas} with $d=H(\epsilon)$, $x_0=0$ and 
     $b_u=B_u$ is a realization of $H$.
     The transformation $R \mapsto \Sigma_R$ preserves 
     span-reachability, observability, minimality and isomorphism.
     For $\Sigma_R$ affine- and span-reachability are equivalent, as its initial state is zero. 
     %and i
     %and the 
     %Mtransformation $\Sigma \mapsto R \mapsto R_{R_{\Sigma}}$
     %results in an SLS which is the shifted version of $\Sigma$: 
     %$x_0,b_u,d$ are replaced by $0$,
     %$b_u+A_u x_0$, $d=Cx_0+d$ respectively. 
    %$\Sigma$ ($\Sigma_R)$ is %respectively affine-reachable, observable, minimal, iff $R_{\Sigma}$ ($R$
    %is 
    %The transformation $\Sigma %\leftrightarrow R_{\Sigma}$ 
%
 %   preserves span-reachability, observability, minimality, and isomorphism.
     %We call SLS the \emph{SLS associated with $\Sigma$}.
     %for all $u(0),\ldots, u(t) \in \mathbb{U}$, $t \ge 0$, where 
     %$u_{s+1} \cdots u(t-1)$ is considered to be the empty sequence $\epsilon$ for $s+1 > t-1$. 
     %In turn, the \eqref{sas:io:eq1} is equivalent to 
     %$H_u(w)=CF_wb_u$ for all $w \in \mathbb{U}^{*}$, $u \in \mathbb{U}$, which 
     %is exactly equivalent to $\Sigma_{lin}$ being a SLS realization of $\Psi_H$.  
     
     %nd hence $\Sigma$ is a minimal realization of $H$, if and only if $R_\Sigma$ is a minimal realization of $\Psi_H$.
     %Moreover, span-reachability of $\Sigma$ from the set of initial states $B=\{0\}$ is 
     %equivalent to span-reachability of $R_\Sigma$, and observability of $\Sigma$ is
     %Mequivalent to observability of $R_\Sigma$. A matrix $T$ defines a linear morphism
     %between two SASs, if and only if it is a morphism between the associated SLSs. 
     From Theorem \ref{theo:real1} it thus follows \cite{Son:Real}
     that  \textbf{(1)} $H$ has a realization by a SAS if and only 
     if $\Psi_H$ has a finite Hankel-rank, \textbf{(2)} a SAS is a minimal realization of $H$ 
     if and only if it is affine-reachable and observable, \textbf{(3)} all minimal SAS realizations of $H$ are isomorphic. 

     We  call \eqref{eq1:sas} \emph{globally uniformly exponentially stable (GUES) with decay rate $\beta$} if 
     the associated SLS is GUES. 
      This definition does not require \eqref{eq1:sas} to be stable in the classical sense as an equilibrium point may not exist. However, if $\mathbb{U}$ is finite and \eqref{eq1:sas} is GUES, then by \cite{ATHANASOPOULOS2016158,Rossa2023} it has an exponentially attractive minimal invariant set, which consists of the equilibrium point if the latter exists.
       %If $x_0 = A_u x_0 + b_u$ for all $u$, then $\{x_0\}$ is invariant and globally attractive \cite{ATHANASOPOULOS2016158,Rossa2023}.}
     %\textcolor{black}{This definition does not imply that \eqref{eq1:sas} is GUES for some equilbirum point in the classical sense. In fact, an equilibrium point of \eqref{eq1:sas} may not exist. However, \cite{ATHANASOPOULOS2016158,Rossa2023} implies that if $\mathbb{U}$ is finite, and \eqref{eq1:sas} is GUES, then there exists an exponential attractive minimal invariant set, in particular, if there is an equilibrium $x_0$ such that
    % $x_0=A_ux_0+b_u$ for all $u \in \mathbb{U}$, then $\{x_0\}$ is a minimal invariant set and all solutions of  \eqref{eq1:sas} converge to $x_0$ 
     %}
     
     \textcolor{black}{For the response 
     maps of SASs, the counterpart  of the decay property is the  forgetting property. }
      %We will call the response map  $H:\mathbb{U}^{+} \rightarrow \mathbb{R}^p$ 
      We say that $H$ is \emph{uniformly forgetting (UF)}
      % if with
      \textcolor{black}{if $\sup_{v \in \mathbb{U}^{*},u \in \mathbb{U}}\|H(uv)-H(v)\|_2 < +\infty$ and}
    %  $\sup_{v \in \mathbb{U}^{+},|v|=t} \sup_{w \in \mathbb{U}^{*}} \|H(wv)-H(v)\| < \infty$ for all  $t \ge 1$, and }
      %\textcolor{black}{
      %$\sup_{v \in \mathbb{U}^{+},|v|=t} \sup_{w \in \mathbb{U}^{*}} \|H(wv)-H(v)\| < \infty$, and } 
       \begin{equation}
       \label{def:uf}
       %\begin{split}
      %    & \forall t \ge 0: 
      %    a_t\!\!=\!\!\sup_{v,w \in \mathbb{U}^{*},|v|=t}\|H(wv)-H(v)\|_2 < \infty \\
       %    & 
      % \foralœl t \ge 0: a_t < +\infty \mbox{ and } \lim_{t \rightarrow \infty} a_t = 0 
      %\begin{aligned}
     %  & \forall t \ge 0:    \sup_{v,w \in \mathbb{U}^{*},|v|=t}\|H(wv)-H(v)\|_2 < \infty \mbox{ and }  \\
       \lim_{t \rightarrow \infty} \sup_{\textcolor{black}{v \in \mathbb{U}^{+}, |v|=t}, w \in \mathbb{U}^{*}} \|H(wv)-H(v)\|_{\textcolor{black}{2}}=0\textcolor{black}{,}
       %\end{aligned}
       \end{equation}
       and we say that $H$ is \emph{uniformly exponentially forgetting (UEF)} with rate $\beta \in (0,1)$, if there exists $C_{f} > 0$, %$\beta \in (0,1)$
       \begin{equation}
       \label{def:uef}
       \forall \textcolor{black}{v \in \mathbb{U}^{*}}:  \sup_{w \in \mathbb{U}^{*}} \|H(wv)-H(v)\|_{\textcolor{black}{2}} \le C_f \beta^{|v|}\textcolor{black}{.} 
       \end{equation}
       %We call $\beta$ \emph{forgetting rate}.
%       Clearly, UEF $\implies$ UF.
  %\begin{Remark}
  %\textit{Remark: Intuition (for exponential) forgetting.}
      \textcolor{black}{(Exponential)} forgetting  means that the influence of earlier inputs on the current output \textcolor{black}{(exponentially)} fades with time, and it 
      \textcolor{black}{relates to decaying properties of $\Psi_H$ as follows.}
       %if two input sequences share the same last $k$ symbols but have different prefixes, then $\|H(wv)-H(v)\|$ decays (exponentially) as $|v|=k$ grows. 
      %This formalizes that stable state-affine dynamics depend mainly on recent inputs.
  %\end{Remark}
  %Forgetting property of $H$ 
  \begin{Lemma}
      \label{UD:UF}
      \textcolor{black}{If $\Psi_H$ is UED with rate $\beta$ if and only if $H$ is  UEF with rate $\beta$}. Moreover, if $H$ is UF, then $\Psi_H$ is UD.
      \end{Lemma}
    \begin{proof}
        %The lemma follows by noticing that
         % $H_u(v)=H(uv)-H(v)$ and that 
        %$H(wv)-H(v)=\sum_{i=1}^{|w|} H_{w[i]}(w[i+1] \cdots w[|w|]v)$, 
         % for all $w \in \mathbb{U}^{*}$, $v \in \mathbb{U}^{+}$. 
             If $\Psi_H$ is UED with rate $\beta$, then 
             %\textcolor{black}{for all $v \in \mathbb{U}^{*}$,}
             \textcolor{black}{by using
             \( \|H(wv)-H(v)\|_2 \le  \textcolor{black}{\sum_{i=1}^{|w|}} \|H_{w[i]}(w[i+1] \cdots w[|w|]v)\|_2 \), 
             \textcolor{black}{$v \in \mathbb{U}^{*}$}, 
             %where $w_{|w|+1:|w|}=\epsilon$, 
             %$w_{i+1:|w|}=w[i+1]\cdots w[|w|]$, $i < |w|$, 
             and
              bounding the terms in the sum by $C_{ed}\beta^{|v|+|w|-i}$,}
             %and setting $C_f=C_{\ed}/(1-\beta)$
             we get 
             \eqref{def:uef} \textcolor{black}{with $C_f=C_{ed}/(1-\beta)$. 
             Conversely, \eqref{def:uef} implies $\|H_{u}(v)\|_2=\|H(uv)-H(v)\|_2 \le C_f\beta^{|v|}$, $v \in \mathbb{U}^{*}$, i.e., $\Psi_H$ is UED with rate $\beta$.}
            % 
            % & \sum_{i=1}^{|w|}C\beta^{|v|+|w|-i} \le \beta^{|v|} C/(1-\beta) 
            %\end{split}
          %\]
           %if $H$ is UEF, then $\|H_u(w)\|_2=\|H(uw)-H(w)\|_2 \le C\beta^{|w|}$ for all $u \in \mathbb  {U}$, $w,v \in \mathbb{U}^{*}$.
            If $H$ is UF, then
            \textcolor{black}{from \eqref{def:hu} it follows that}
                %\textcolor{black}{$\sup_{u \in \mathbb {U}, v \in \mathbb{U}^*} \|H_u(v)\|_2 \le \sup_{w,v \in \mathbb{U}^*} \|H(wv)-H(v)\|_2 < +\infty$
                %and 
            %$H_u(v)=H(uv)-H(v)$, 
            %\[ 
             %   \lim_{t \rightarrow \infty} 
                %le \lim_{t %\rightarrow \infty} 
            $\sup_{v \in \mathbb{U}^{*}, |v|=t} \|H_u(v)\|_2 \le
            \sup_{\substack{v \in \textcolor{black}{\mathbb{U}^{+}}, |v|=t \\ w \in \mathbb{U}^{*}}}
                \|H(wv)-H(v)\|_2$
                \textcolor{black}{for $t \ge 1$. Hence,}
                \textcolor{black}{by} taking the limit $t \rightarrow \infty$ we get \eqref{def:ud}, \textcolor{black}{and 
                 $\sup_{\substack{u \in \mathbb{U} \\ v \in \mathbb{U}^{*}}} \|H_u(v)\|_2\!=\! \sup_{\substack{u \in \mathbb{U} \\ v \in \mathbb{U}^{*}}}\! \! \|H(uv)-H(v)\|_2<+\infty$.}
                 %for all $u \in \mathbb{U}$, $w \in \mathbb{U}^{*}$, and hence $\lim_{t \rightarrow \infty} \sup_{u \in \mathbb{U}, w \in \mathbb{U}^{*}, |w|=t} \|H_u(w)\|_2=0$, i.e., $\Psi_H$ is UD.
    \end{proof}
   From Theorem \ref{main:theo} we get thus  the following.
   \begin{Corollary}
   \label{col:sas}
      The response map 
        $H$ can be realized by a GUES SAS 
        if and only if $H$ is UF  and 
        $\Psi_H$ has a finite Hankel-rank.
        If $\Psi_H$ has a finite Hankel-rank, $H$ is UF if and only if $H$ is UEF.
        %$\dim W_H=n < +\infty$.
        %the maximal rank of its Hankel-matrix is bounded, i.e.,
        %$\sup_{(\alpha,\beta) \mbox{$n$-selection}, n \ge 1 } \Rank~ \hankred^{H}_{\alpha,\beta} < +\infty$.
	   If $H$ is UEF with \textcolor{red}{rate} $\beta$, then all minimal SAS realizations of $H$ are GUES with rate $\beta$. 
         \end{Corollary}
         \begin{color}{black}
       \begin{Remark}[Echo-state, fading memory]
        \label{rem:echo_fading}
      The UF property is related to the \emph{fading-memory (FM)} property for maps on infinite input sequences \cite{Ortega3, Ortega4}. If $H$ is UF and realizable by a SAS, then any minimal realization \eqref{eq1:sas} of $H$ is GUES; thus, by \cite[Theorem 1]{Ortega4}, it has the echo-state property and induces the infinite-sequence response map
      %\begin{align*}
      \( H_{ext}(\underline{u}) = Cb_{\underline{u}[0]} + \sum_{k=-1}^{-\infty} CA_{\underline{u}[0]\cdots \underline{u}[-k+1]}b_{\underline{u}[-k]} = 
      \lim_{t \to -\infty} H(\underline{u}[-t] \cdots \underline{u}[0]) \),
      %\end{align*}
      which is well-defined and has the FM property. We conjecture that FM maps on infinite sequences are those realizable by GUES SASs, and that they induce UF maps on finite sequences.
      \end{Remark}  
    \end{color}
  %\paragraph{Bounded state-affine systems.}
   %for which for which $h(x,u)=C_0x$ is linear in $x$ and does not depend on $u$.
   %$C_i=0,d_i=0$ for all $i=1,\ldots, n_p$ and $d_0=0$.  
    Among SASs $\psi$-bounded SASs are of particular interest, as they include \textcolor{black}{such popular classes as} bilinear systems, and they are computationally effective.
  By \cite{Son:Real}, $H$ can be realized by a $\psi$-bounded SAS if and only if $H$ is bounded of type $\{\psi_{q}\}_{q=0}^{n_p}$ according to \textcolor{black}{\cite[Definition 1.15]{Son:Real}}, and the Hankel-rank of $\Psi_H$ equals the rank of a suitable Hankel matrix. \textcolor{black}{Minimization of  $\psi$-bounded SASs results in $\psi$-bounded SASs, which} are minimal SASs realizations. Thus:
%   SASs which are $\psi$-bounded  are of interest, as they include \textcolor{black}{such popular classes as} bilinear systems, and they are computationally effective.
%  By \cite{Son:Real}, $H$ can be realized by a $\psi$-bounded SAS if and only if $H$ is bounded of type $\{\psi_{q}\}_{q=0}^{n_p}$ as in \textcolor{black}{\cite[Definition 1.15]{Son:Real}}, and the Hankel-rank of $\Psi_H$ equals the rank of a suitable Hankel matrix. \textcolor{black}{Minimization of  $\psi$-bounded SASs results in minimal SASs which are also $\psi$-bounded SASs} % which} are minimal SASs realizations. Thus:
  %and any $\psi$-bounded SAS can be minimized. Thus:
   %,i.e.,  
   %there exists maps $S_{H,q}:Q^{*} \rightarrow \mathcal{Y}$, $q \in Q$,
   %$Q=\{0,1,\ldots, n_p\}$, 
  %called \emph{generating series}
  %of $H$, such that $H(u(0)\cdots u(t))$ is 
  %\begin{equation*}
  %\begin{split}
  %    &  \\
  %\sum_{i=0}^{t-1} \sum_{q_i,\ldots, q_{t-1}} S_{H,q_i}(q_{i+1} \cdots q_{t-1}) \prod_{s=i}^{t-1} \psi_{q_s}(u(s)).
  %\end{split}
  %\end{equation*}
  %We can then associate with $H$ the family $\Psi_{H,b}=\{\tilde{S}_{H,q} \mid q \in Q\}$ of response maps defined on a finite set of inputs ($Q$).
  %Then $H$ has a realization by a $\psi$-bounded SAS \eqref{eq1:sas}-\eqref{eq:sas_psi} if and only if $\Psi_{H,b}$ has a realization by a SLS.
  % $R_{\Sigma,b}=(n,\{F_i\}_{q \in Q},C,\{b_q \mid q \in Q\})$. Note that the latter SLS $R_{\Sigma,b}$ gives rise to a SLS $R_{\Sigma}=(n,\{F_u=\sum_{i=0}^{n_p} F_i\psi_i(u)\}_{u \in \mathbb{U}},\{\sum_{i=0}^{n_p} b_i\psi_i(u) \mid u \in \mathbb{U}\},C)$ which is a realization of $\Psi_H$.
  %Moreover, there is a one-to-one correspondence between $\psi$-bounded SAS realizations of $H$ and SLS realizations of $\Psi_{H,b}$. This correspondence preserves span-reachability, observability, minimality, isomorphism.
  %By \cite[Section I.C]{Son:Real}, if $H$ is bounded of type $\{\psi_{q}\}_{q \in Q}$, then $H$ has a realization by a SAS if and only if it has a realization by a $\psi$-bounded SAS. 
  %The finite Hankel-rank condition of $\Psi_H$ is in fact equivalent to the finite Hankel-rank condition of $\Psi_{H,b}$, and hence to the existence of a $\psi$-bounded SAS realization of $H$.
  \begin{Corollary}
   The statements of Corollary \ref{col:sas} remain true if  SAS is replaced by $\psi$-bounded SAS and the finite Hankel-rank of $\Psi_H$ is replaced by the finite rank of the Hankel-matrix defined in \cite[Definition 2.7]{Son:Real}. 
  \end{Corollary}
  % which is equivalent to the existence of a $\psi$-bounded SAS realization of $H$.

%
%
%$\Psi_H$ has a realization by a SLS, i.e., if
%
% and $\Psi_H$ has a finite Hankel-rank. Moreover, if $H$ is exponentially forgetting with rate $\beta$, then all minimal $\psi$-bounded SAS realizations of $H$ are uniformly exponentially stable with rate $\beta$.


 

%%Likewise, we will say that the response map $H:\mathbb{U}^{+} \rightarrow \mathbb{Y}$
%%is \emph{$\psi$-bounded}, if 
%%
%%
%%Note that for any SAS $\Sigma$ there exists functions $\psi=\{\psi_q\}_{q=1}^{n_p}$ for some integer $n_p$, such that $\Sigma$ is $\psi$-bounded. In particular, if
%%$F_u,b_u,C_u,d_u$ are polynomials functions of $u$, then 
%%$\psi_q$ can be chosen to be polynomial.
%%
%%Likewise, if $H$ is polynomial in inputs with a bounded degree, i.e., 
%%there exists $d$ such that for all $w \in \mathbb{U}^{+}$,  $H(w)$ is polynomial 
%%in $w[1],\ldots, w[|w|]$ of degree $d$ in each $w[i]$, $i=1,\ldots |w|$, then there exists
%%polynomials $\psi=\{\psi_i\}_{i=1}^{n_p}$ such that $H$ is $\psi$-bounded.
%%
%%That is, the assumption that the response map or the SAS is $\psi$-bounded is not overly restrictive.
%%It is easy to see \cite[Section C]{Son:Real} that $H$ has a realization by a $\psi$-bounded SAS, if $H$ is $\psi$-bounded. 
%%
%%We say that $\Sigma$ is a \emph{minimal $\psi$-bounded realization} of $H$, if $\Sigma$ has the minimal dimension among all $\psi$-bounded SAS realizations of $H$. 
%%
%%
%%It is well-known that realization theory of bounded SAS can be reduced to that of
%%recognizable formal power series as follows. 
%%\begin{Lemma}[based on \cite{Son:Real}]
%%  $H$ can be realized by a $\psi$-bounded SAS only if $H$ is $\psi$ bounded. 
%%  A $\psi$-bounded response map $H$ can be realized by a $\psi$-bounded SAS if and only if 
%%  $\Psi_H=\{\tilde{S}_{H,q} \mid q \in Q\}$, where 
%%  \[ \tilde{S}_{H,q}(w) = \begin{bmatrix} (S_{H,q}(w1))^T & \cdots  & (S_{H,q}(w\QNUM)))^T \end{bmatrix}^T \]
%%  has a linear representation (its elements are recognizable
%%  formal power series). Moreover, there is a one-to-on correspondence
%%  between $\psi$-bounded SAS realizations of $H$ satisfying \eqref{eq1.2}
%%  \footnote{if $\Sigma$ is a realization of $H$, then necessarily $d_q=S_{H,q}(\epsilon)$}
%%  and
%%  SLS $R=(n,\{F_q\}_{q \in Q},\{b_q \mid \in Q\},\begin{bmatrix} C_1^T & \cdots & C_{\QNUM}^T \end{bmatrix}^T))$ realizations of $\Psi_H$. This correspondence preserves span-reachability, observability, minimality, isomorphism. 
%%\end{Lemma}
%%\begin{Corollary}[\cite{Son:Real}]
%%  A response map $H$ can be realized by a $\psi$-bounded SAS, if and only if 
%%  $H$ is $\psi$-bounded and 
%%  $\sup_{(\alpha,\beta) \mbox{$n$-selection}, n \in \mathbb{N} } \Rank~  \hankredu^{\Psi}=n < +\infty$. 
%%  A $\psi$-bounded SAS is a minimal realization of $H$ if it is span-reachable
%%  and observable. All minimal $\psi$-bounded SAS realizations of $H$ are isomorphic. 
%%  In particular, minimal $\psi$-bounded SAS realizations of $H$ are also minimal
%%  among SAS realizations of $H$.
%%\end{Corollary}
%%We call 
%%$\sup_{(\alpha,\beta) \mbox{$n$-selection}, n \in \mathbb{N} } \Rank~  \hankredu^{\Psi}=n < +\infty$
%%the \emph{finite Hankel rank} condition. 
% For $\psi$-bounded response maps and SASs the discussion above can be made computationally
% effective. 
% Let  $\psi=\{\psi\}_{i=1}^{n_p}$ be a family of linearly independent functions. 
  % \begin{Notation}
  %  Let $Q=\{1,\ldots,n_p\}$ and for every $w=q_1\cdots q_k \in Q^{*}$ and every 
  %  $\underline{u}=u_1\cdots  u_t \in \mathbb{U}^{*}$
  %  let $\psi_{w}(\underline{u})$ be defined as follows:
  %  \[ \psi_w(\underline{u})=\prid_{i=1}^{t} \psi_{q_i}(u_i) \]
  % \end{Notation}
%Let be $H$ be a  \emph{$\psi$-bounded response map}.
%%In addition, observability and span-reachability of $\psi$-bounded SASs can be checked by checking
%%the rank of the following finite matrices.
%%\begin{Theorem}[\cite{Son:Real}]
%%  Let $\Sigma$ be a $\psi$-bounded SAS realization of $H$.
%%  \begin{itemize}
%%  \item
%%      The SAS $\Sigma$ is minimal if and only if it is span-reachable
%%      and observable. 
%%   \item  
%%      The SSA $\Sigma$ is span-reachable, if and only if
%%      \[ \dim \SPAN\{ F_wb_q \mid w \in Q^{*},q \in Q, |w| < n \}=. \]
%%      Likewise, $\Sigma$ observable, if and only if 
%%      \[ \dim \SPAN\{ F_w^TC_q^Tv \mid w \in Q^{*},q \in Q, |w| < n, v \in \mathbb{R}^p \}=n. \]
%%\item  
%%   The SAS $\Sigma$ can be transformed to a span-reachable and observable $\psi$-bounded SAS
%%   realization $H$ by a minimization algorihm 
%%\end{itemize}
%%\end{Theorem}
%Theorem \ref{main:theo} leads to the following corollary.
%\begin{Corollary}[Stable bounded SSA]
%   A response map $H$ has a realization by a $\psi$-bounded SAS which is GUES if and only if
%   $H$ is $\psi$-bounded, UF, and  it satisfies the finite Hankel-rank condition.
   %$n=\sup_{(\alpha,\beta) \mbox{$n$-selection}, n \ge 1 } \Rank~ \hankredu^{\Psi} < +\infty$.
%   If $H$ is $\psi$-bounded, and UF and it satisfies the finite Hankel-rank condition, then 
%   all minimal $\psi$-bounded SAS realizations are GUES. If, in addition, 
%   $H$ is UEF with rate $\beta$, then all minimal $\psi$-bounded SAS realizations of
%   $H$ are GEUS with rate $\beta$.
%\end{Corollary}

\subsection{LPV-SSA/LSS} %Special case: LPV-SSA/LSS}
\textcolor{black}{We} apply Theorem \ref{main:theo} to LPV-SSA/LSSs.
Their 
response maps %of LPV-SSA/LSS are maps 
are of the form $H: (\mathbb{P} \times \mathbb{R}^m)^+ \to \mathbb{R}^{\textcolor{black}{n_y}}$.
Following \cite{Petreczky2016}, we  say that $H$ has an \textcolor{black}{\emph{impulse response representation (IRR)}} if there exists a map $S: \textcolor{black}{\bigcup_{k=2}^{\infty} Q^{k}} \to \mathbb{R}^{\textcolor{black}{n_y} \times m}$, 
where $Q=\{1,\ldots,n_p\}$, such that 
\textcolor{black}{for any $\underline{p} \in \mathbb{P}^{\infty}$, $\underline{v} \in (\mathbb{R}^m)^{\infty}$}, %$t \ge 1$, %where 
%and $h_j:\mathbb{P}^{*} \to \mathbb{R}^p$ such
\begin{align*}
 & H((\underline{p}[1], \underline{v}[1])\cdots (\underline{p}[t], \underline{v}[t])) = \sum_{\tau=1}^{\textcolor{black}{t-1}} (h \diamond \underline{p})(t,\tau)\underline{v}[\tau], \quad \textcolor{black}{t \ge 1} \\
%\end{equation*}
%\textcolor{black}{for $\tau \in[1, t-1]$}
%\begin{equation*}
   & (h \diamond \underline{p})(t,\tau)=\sum_{s \in Q^{t-\tau+1}} S(s)\underline{p}_s(t), \quad \textcolor{black}{\tau \in [1,t-1]}
%\vspace{-2pt}
\end{align*}
\textcolor{black}{and $\underline{p}_s(t)=\prod_{i=1}^{|s|} \textcolor{black}{(\underline{p}[t-|s|+i])_{s[i]}}$\textcolor{black}{; recall that $(\underline{p}[k])_r$ is the $r$th component of $\underline{p}[k]$}.}
%\textcolor{black}{where
  %and $\underline{p}_s(t)=\prod_{i=1}^{k} \underline{p}_{s_i}[t-k+i]$ for $s=s_1 \cdots s_{k} \in Q^{+}$.}
 % \in (\mathbb{P} \times \mathbb{R}^m)^+$:
   %\prod_{i=\tau}^{t} p_{q_i}(i)
    %\substack{s \in Q^{t-\tau-1} \\ q_\tau,\, q_t}} 
  %(w_{q_t s q_\tau}\diamond p)(t,\tau) \\
  %& (w_{\epsilon}\diamond p)(t,\tau)=
  %\begin{cases}
  %1, & t=\tau-1,\\
  %0, & t\leq \tau-1,
  %\end{cases} \\
  %& 
  %\text{if } s=q_1 \cdots q_{t-\tau-1}  \text{ with } q_i \in Q,\ \text{then} 
  %(w_{q_{\tau}q_{\tau+1}\cdots q_t}\diamond p)(t,\tau)=
  %\prod_{i=1}^{t-\tau-1} p_{q_i}(\tau+i).
  We call $h \diamond \underline{p}$ the \textcolor{black}{\emph{impulse response} of $H$ (IRR)}  for the scheduling sequence $\underline{p}$\textcolor{black}{, 
  and we call the values of $S$ \emph{sub-Markov parameters.}}
  \textcolor{black}{By \cite{Petreczky2016}, 
  $\Sigma$ from \eqref{eq1:lpv} is a realization of $H$
  if and only if $H$ has an IRR and $S(q_1vq_2)\!=\!C_{q_2}A_vB_{q_1}$, 
  for all $q_1,q_2 \in Q$, $v \in Q^{*}$.}
  If $\mathbb{P}=\{e_1,\ldots,e_{n_p}\}$ (switched case) and \textcolor{black}{$q \in Q^{\infty}$ is such that}
  $\underline{p}[s]=e_{\underline{q}[s]}$, then
  $(h \diamond \underline{p})(t,\tau)=S(\underline{q}[\tau]\cdots \underline{q}[t])$, i.e., the sub-Markov parameters are the impulse responses.
  For general $\mathbb{P}\subseteq\Delta$, $(h \diamond \underline{p})(t,\tau)$ is a convex combination of these parameters.
  % $p(\tau)=e_{q_\tau}$, and $p(s)=e_{q_s}$ for $s=\tau+1,\ldots t-1$. In this case, the IRR condition is equivalent to the existence of sub-Markov parameters $S_{q,q_0}: Q^* \to \mathbb{R}^p$ such that $H((p(0), v(0)), \ldots, (p(t), v(t))) = \sum_{s=0}^{t-1} S_{q_t,q_s}(q_{s+1} \cdots q_{t-1})v(s)$, which
  %is the usual convolution of the input sequence $v$ with the impulse response $h$. In general,
  %the IRR condition is equivalent to the existence of sub-Markov parameters $S_{q,q_0}: Q^* \to \mathbb{R}^p$ such th
  Following \cite{Petreczky2016} we define the \textcolor{black}{maps} $M_{q,\textcolor{black}{l}}: Q^* \to \mathbb{R}^{n_p \textcolor{black}{n_y}}$,
  $q \in Q$, $l=1,\ldots,m$ \textcolor{black}{as} 
   \[
   M_{q,l}(w) = \begin{bmatrix} (S_{\cdot,\textcolor{black}{l}}(qw1))^T, & \ldots, & (S_{\cdot,\textcolor{black}{l}}(qwn_p))^T \end{bmatrix}^T.
  \]
  %as \textcolor{black}{follows:}
  \textcolor{black}{
  for all 
  %$q \in Q$, $\textcolor{black}{l=1,\ldots, m}$ 
  $w \in Q^{*}$, 
  where $S_{\cdot,\textcolor{black}{l}}(s)$ is the \textcolor{black}{$l$}th column of $S(s)$,}
  and we consider the following family of response maps:
  \textcolor{black}{$\Psi_H = \{M_{q,l}\}_{j=(q,l) \in J}$ with $J=Q \times \{1,\ldots,m\}$.
   For a LPV-SSA/LSS $\Sigma$ from \eqref{eq1:lpv}, for any $(q,l) \in J$
   let $\tilde{B}_{(q,l)}$ be the $l$th column of $B_q$, and 
   associate with $\Sigma$ the SLS
  }
  %We associate with any LPV-SSA/LSS the SLS 
  \[ 
  \begin{split}
  &R_{\Sigma}=(\{A_q\}_{q \in Q},\textcolor{black}{\{\tilde{B}_{(q,l)}\}_{(q,l) \in J}},\begin{bmatrix} C_1^T, & \ldots & C_{n_p}^T \end{bmatrix}^T). \\
  %\widetilde{B},\widetilde{C}) \\
  %& \widetilde{C}=\begin{bmatrix} C_1^T & \ldots & C_{n_p}^T \end{bmatrix}^T
  % \quad \widetilde{B}=
  \end{split}
  \]
  % where $A(p)=\sum_{i=1}^{n_p} A_i p_i$, $B(p)=\sum_{i=1}^{n_p} B_i p_i$.
  %A LPV-SSA/LSS 
  \textcolor{black}{Then} $\Sigma$ is a realization of $H$ if and only if \textcolor{black}{the SLS $R_{\Sigma}$ is a realization of $\Psi_H$} \cite{Petreczky2016}. The correspondence between $\Sigma$ and $R_{\Sigma}$ is a one-to-one \textcolor{black}{map between LPV-SSA/LSS realizations of $H$ and SLS realizations of $\Psi_H$, and it} preserves span-reachability, observability, minimality, and isomorphism \cite{Petreczky2016}.
  \textcolor{black}{Hence,} $H$ has a realization by an LPV-SSA/LSS if and only if $H$ is IRR and $\Psi_H$ has a  finite Hankel-rank, $\Sigma$ is a minimal realization of $H$ if and only if it is span-reachable, observable, and any two minimal LPV-SSA/LSS realizations of $H$ are isomorphic \cite{Petreczky2016}. 
  %Any LPV-SSA realization of $H$ can be transformed to a minimal LPV-SSA realization of $H$. The dimension of a minimal LPV-SSA realization of $H$ equals the finite Hankel-rank of $H$. 
  %Conversely, we can associate with any state-linear system $R$ an LPV system $\Sigma_R$ such that the correspondence $R \to \Sigma_R$ is the inverse of $\Sigma \to R_{\Sigma}$.

  The LPV-SSA/LSS  \eqref{eq1:lpv} is called \emph{globally uniformly exponentially stable (GUES)} if there exist constants \textcolor{red}{$C \ge 1$}, $\beta \in (0,1)$ such that 
  \textcolor{black}
  {\(
  \|A(\underline{p}[k]) \cdots A(\underline{p}[1])\|_2 < C \beta^k
  \),
  for any  $\underline{p} \in \mathbb{P}^k$, $k >0$, 
  This is equivalent to \eqref{eq1:lpv} %the LPV-SSA/LSS 
  being GUES at the zero equilibrium point for $v=0$, 
  and can be characterized by}
  the joint spectral radius of $\{A(p)\}_{p \in \mathbb{P}}$ being smaller than $1$. \textcolor{black}{As} $\{A_i\}_{i=1}^{n_p} \subseteq \{A(p)\}_{p \in \mathbb{P}} \subseteq Conv(\{A_i\}_{i=1}^{n_p})$, \textcolor{black}{the latter} is equivalent to SLS $R_{\Sigma}$ being GUES \cite{RJungers}.
  %
  %is equivalent to the GUES property of the corresponding state-linear system $R_{\Sigma}$. This follows from the fact that $A(p)$ is a convex combination of $\{A_i\}_{i=1}^{n_p}$ and \cite{RJungers},the GUES property of $\Sigma$ is equivalent to the joint spectral radius of $\{A(p)\}_{p \in \mathbb{P}}$ being smaller than $1$.

    
  %if $\Psi_H$ is UD, UAS and UED respectively.
  %Intuitively, this means that the Markov parameters 
  %$\{S_{q,q_0}(q_1\cdots q_t)\}_{t=0}^{\infty}$
  %converge to zero as $t \rightarrow +\infty$ (IRR-UD), the convergence is exponential (IRR-UED), respectively the Markov-parameters are absolutely summable 
  %(IRR-DD), and in all cases the convergence is uniform over all the choices  $q,q_0,q_1,\ldots \in Q$.
  %These concepts can be related to properies of the impulse response $h \diamond \underline{p}$ as follows. 
  %Let us say that the impulse response $H$ is
  The IRR of $H$ is called
  %a IRR which is respectively
  %\begin{itemize}
  \emph{uniformly decaying (IRR-UD)}, if 
   \begin{align*}
   & \lim_{N \rightarrow \infty} \sup_{\underline{p} \in \mathbb{P}^{\infty}, \tau \in \mathbb{N},\tau \ge 1}\|(h \diamond \underline{p})(N+\tau,\tau)\|_2=0, 
   \end{align*}
  and \emph{uniformly exponentially decaying (IRR-UED) with rate $\beta \in (0,1)$} if 
  %µand   with rate $\beta \in (0,1)$, if 
  for some $C > 0$ \textcolor{black}{, for all $N \ge 1$,} %and $\beta \in (0,1)$,
  \begin{align*}
   \sup_{\underline{p} \in \mathbb{P}^{\infty}, \tau \in \mathbb{N}, \tau \ge 1} \|(h \diamond \underline{p})(N+\tau,\tau)\|_2 < C\beta^{\textcolor{black}{N-1}}.
   \end{align*}
  
  %\emph{uniformly absolutely summable (IRR-UAS)}
   %\emph{uniformly decaying}, if 
   %\emph{uniformly exponentially decaying}, if
   %and it is \emph{uniformly absolutely summable} 
  % if  $\lim_{n \rightarrow +\infty} b_n=0$ where for any $n \in \mathbb{N}$,
  % \begin{equation} 
  %  \label{eq:IRR-UAS}
  % \begin{split}
  % b_n=\sup_{\underline{p} \in \mathbb{P}^{\infty}, N \in \mathbb{N}}\sum_{r=0}^{N-1} \|(h \diamond \underline{p})(N+n,r)\|_2 < +\infty,
  %  % \sup_{\underline{p} \in \mathbb{P}^{\infty}} \sum_{r=n}^{N+n-1} \|(h \diamond \underline{p})(N+n,r)\|_2=0
  %\end{split}
  %\end{equation}
  Intuitively, IRR-UD (IRR-UED) means that the impulse response $h \diamond \underline{p}$ decays (exponentially) 
  to zero, and it is equivalent to the sub-Markov parameters decaying (exponentially).
  %and IRR-UAS means that the impulse responses are absolutely summable, and all thise convergence properties hold uniformly over the scheduling signals. 
  %\rightarrow +\infty$ for any choice of $\underline{p} \in \mathbb{P}^{\infty}$, IRR-UED means that the convergence is exponential, and IRR-UAS means that the impulse response is absolutely summable. In all cases, the convergence is uniform over all choices of $\underline{p} \in \mathbb{P}^{\infty}$.
  %for any choice of $\underline{p} \in \mathbb{P}^{\infty}$, the impulse response $h \diamond \underline{p}$ is decaying (resp. exponentially decaying, resp. absolutely summable) and the convergence is uniform over all choices of $\underline{p} \in \mathbb{P}^{\infty}$.
  \begin{Lemma}
  \label{lem:lpv:IRR-UD}
  $H$ is IRR-UD (resp. IRR-UED \textcolor{black}{with rate $\beta$}) if and only if $\Psi_H$ is UD 
  (resp. UED \textcolor{black}{with rate $\beta$}). % resp. UAS).
  %%\begin{enumerate}
  %%\item $H$ is IRR-UD 
  %%
  %%$\iff$  $h \diamond \underline{p}$ is uniformly decaying for all $\underline{p} \in \mathbb{P}^{\infty}$, i.e., 
  %%\item $H$ is IRR-UED $\iff$  the impulse response $h \diamond \underline{p}$ is uniformly exponentially decaying for all $\underline{p} \in \mathbb{P}^{\infty}$, i.e., there exists $\beta \in (0,1)$ such that $\sup_{\underline{p} \in \mathbb{P}^{\infty}} \|(h \diamond \underline{p})(N+\tau,\tau)\|_2 < C\beta^N$ for some $C > 0$.
  %%\item $H$ is IRR-UAS $\iff$  the impulse response $h \diamond \underline{p}$ is uniformly absolutely summable for all $\underline{p} \in \mathbb{P}^{\infty}$, i.e., $\sup_{\underline{p} \in \mathbb{P}^{\infty}, N \in \mathbb{N}}\sum_{r=0}^{N-1} \|(h \diamond \underline{p})(N,r)\|_2$ is bounded and 
  %M$\lim_{n \rightarrow +\infty} \sup_{\underline{p} \in \mathbb{P}^{\infty}} \sum_{r=n}^{N+n-1} \|(h \diamond \underline{p})(N+n,r)\|_2=0$.
  %%\end{enumerate}
  \end{Lemma}
  %That is, the various decaying properties of the impulse response $h \diamond \underline{p}$ are equivalent to the corresponding decaying properties of the sub-Markov parameters of $H$.\\
  \begin{proof}[Proof of Lemma \ref{lem:lpv:IRR-UD}]
   We show that 
      \begin{equation}
      \label{eq:pf:lpv1}
      \begin{split}
           %&
          %\sup_{\underline{p} \in \mathbb{P}^{\infty}} \|(h \diamond \underline{p})(K,l)\|_2=
          %\sup_{q_l,q_N \in Q, s \in Q^{N-l-1}} \|S_{q_{N},q_{l}}(s)\|_2 \\
          \sup_{\underline{p} \in \mathbb{P}^{\infty}}\!\! %\sum_{r=l}^{K} 
          \|(h \diamond \underline{p})(N+\tau,\tau)\|_2
          \!\! = \!\! 
          \sup_{s \in Q^{N+1}}  \|S(s)\|_2 
          %\!\! \sum_{r=l}^{K} \|S(s_{r:N})\|_2
      \end{split}
      \end{equation}
   for any \textcolor{black}{$N \ge 1,\tau \in \mathbb{N}$}.  %the following holds:
      %where $s_{r:N}=s[r+1] \cdots s[N]$ is the suffix of $s$ starting at position $r+1$.
      %In particular,  $K=\tau$ and $l=\tau$, 
      %$\sup_{\underline{p} \in \mathbb{P}^{\infty}} \|(h \diamond \underline{p})(N-1,\tau)\|_2=\sup_{s \in Q^{N}} 
      The statement of the lemma follows from \eqref{eq:pf:lpv1}
      \textcolor{black}{by noticing that 
       $\frac{1}{\sqrt{m}} \|S(qwq_1)\|_2  \le \sup_{l=1,\ldots,m} \|M_{q,l}(w)\|_2 \le \sqrt{n_p} \sup_{\sigma  \in Q} \|S(q w\sigma)\|_2$, $q,q_1 \in Q,w \in Q^{*}$.}
      %this particular case, 
      %while the statement for UAS follows from the general case of }.   
     % For the latter, \eqref{eq:pf:lpv1} applied to $l=0$ and $K=T-1$, $N=T+n-1$ implies
     % notice that with the notation of \eqref{eq:IRR-UAS}
     % \[ 
     %     b_n=\sup_{T \in \mathbb{N}} \sup_{s \in Q^{T+n+1}}
     %      \sum_{r=0}^{T-1} \|S(s_{r:T+n})\|_2
     % \]
      % In turn, notice that for any 
        %    %$q_0,q_1 \ldots, \in Q$, 
        %       $\bar{q}_i \in Q$, $i \in \mathbb{N}$,  and any $T$      
        %       by setting $q_i=\bar{q}_{T+n-i}$ for $i=1,\ldots, T+n-1$, 
        %       and $s=q_0 \cdots q_{T+n}$,
        %       we have
        %    \[
        %        \begin{split}
        %         \sum_{r=n+1}^{T+n-1} \|M_{q,j}(\bar{q}_{r}\cdots \bar{q}_{1})\|_2= 
        %         \sum_{r=0}^{T-1} \|M_{q_r,j}(s_{r+1:T+n-1})\|_2 
        %       \end{split}    
        %       \]
        %       By uniform convergence of 
        %       $\sum_{r=0}^{\infty} \|M_{q,j}(\bar{q}_{r}\cdots \bar{q}_{1})\|_2$
        %       it follows that $\epsilon > 0$, there exists $K_{\epsilon}$  (independent of $\{\bar{q}_i\}_{i=0}^{\infty},q,j$) such that for any $n > K_{\epsilon}$,
        %       \(
        %          \sum_{r=n+1}^{N+n-1} \|M_{q,j}(\bar{q}_{r}\cdots \bar{q}_{1})\|_2 < \epsilon/m
        %       \), and hence 
        %       \[ \sum_{r=0}^{T-1} \|S(s_{r:T+n})\|_2 \le \sum_{j=1}^{m}\sum_{r=0}^{T-1} \|M_{q_r,j}(s_{r+1:T+n-1})\|_2 <\epsilon
        %       \]
        %       for any $T$ and any
        %       choice of $s \in Q^{T+n+1}$, which implies that $\lim_{n \rightarrow +\infty} b_n=0$.    
        To show \eqref{eq:pf:lpv1}, for any choice of $\underline{q} \in Q^{\infty}$, and any $N$, choose $\underline{p} \in \mathbb{P}^{\infty}$ such that $\underline{p}[i]=e_{\underline{q}[i]}$ \textcolor{black}{for all $i \ge 1$}. Then 
        $(h \diamond \underline{p})(N+\tau,\tau)=S(\textcolor{black}{\underline{q}[\tau] \cdots \underline{q}[N+\tau]})$
        and hence the left-hand side of \eqref{eq:pf:lpv1} is greater than or equal to the right-hand side.
        Conversely, for any choice of $\underline{p} \in \mathbb{P}^{\infty}$
	  by using the fact that $\mathbb{P} \subseteq \Delta$\textcolor{red}{,}  for any $N$, % is the standard simplex, for any $N$, 
        \[ 
          \begin{split}
             & \|(h \diamond \underline{p})(N+\tau,\tau)\|_2  \le \sum_{s \in Q^{N+1}} \|S(s)\|_2 \underline{p}_s(N+\tau) \\
            & \le \sup_{s \in Q^{N+1}} \|S(s)\|_2 \sum_{s \in Q^{N+1}} \underline{p}_s(N+\tau) =
            \sup_{s \in Q^{N+1}} \|S(s)\|_2 \textcolor{black}{,}
          \end{split}
        \]
        where we used that  $\sum_{s \in Q^l} \underline{p}_s(t) = 1$ for any $t \ge l$.
        %and that $\underline{p}_{s_1}(t-|s_2|)\underline{p}_{s_2}(t)=\underline{p}_{s_1 s_2}(t)$.
  %%    Hence,
  %%    %\[
  %%    %  \begin{split}
  %%    %  &  \|h \diamond \underline{p}(N,\tau)\|_2  \le \sup_{s \in Q^{N}} 
  %%    %     \|S(s)\|_2 \sum_{s \in Q^{N}} p_s(N) =
  %%    %     \sup_{s \in Q^{N+1}} \|S(s)\|_2 
  %%    %     %\times  \\ & \times \prod_{i=1}^{N+k} p_{q_i}(i) = 
  %%    %    %\sup_{q_\tau,q_{N} \in Q,s \in Q^{N-\tau-1}} 
  %%    %     %\|S_{q_{N}, q_r}(s)\|_2 
  %%    %  \end{split}
  %%    %\]
  %%    \[
  %%      \begin{split}
  %%      & \sum_{r=l}^{K} \|h \diamond \underline{p}(N-1,r)\|_2 \le \sum_{s \in Q^{N}} 
  %%         \sum_{r=l}^{K}\|S(s_{r:N})\|_2 p_s(N) = \\
  %%         %\prod_{i=1}^{N} p_{q_i}(i) = \\
  %%        &  \le \sup_{s \in Q^{N}} \sum_{r=l}^{K} \|S(s_{r:N})\|_2 
  %%        \!\! \sum_{s \in Q^{N}}\!\! p_s(N) \!\!=\!\! 
  %%         \sup_{s \in Q^{N}} \sum_{r=l}^{K-1} \|S(s_{r:N})\|_2 
  %%      \end{split}
  %%    \]
  \end{proof}
  Theorem \ref{main:theo} applied to $\Psi_H$ and SLSs associated with LPV-SSAs results in the following.
      \begin{Corollary}%[Realizability by stable LPV-SSAs/LSS]
        \label{cor:lpv}
        Let $H$ be a response map \textcolor{black}{with} an IRR. 
        \\
        \textbf{(A)}
         If $\Psi_H$ has a finite Hankel-rank, then
         \textcolor{black}{IRR-UD and IRR-UED are equivalent.}
         %the following
          %properties of $H$ are equivalent: IRR-UD and IRR-UED, 
         \\
        \textbf{(B)}
         $H$ has a realization by a GUES LPV-SSA/LSS if and only  if 
         $\Psi_H$ has a finite Hankel-rank, and 
         \textcolor{black}{IRR-UD holds.}
         %any of the equivalent conditions holds: IRR-UD, IRR-UED.
         \\
        \textbf{(C)}
         $H$ is IRR-UED with the rate $\beta$ if and only if all minimal 
         LPV-SSA/LSS realizations of $H$ are GUES with the rate $\beta$.
       %\end{itemize}
      \end{Corollary}
      That is, the existence of a GUES LPV-SSA/LSS realization is equivalent to the existence of an IRR which is decaying, necessarily exponentially, 
      %is uniformly decaying or uniformly exponentially decaying,
      and the decay rate for IRR is the same as that of any minimal realization.
      %Moreover, the existence of an IRR which is uniformly exponentially decaying with rate $\beta$ is equivalent to the existence of a GUES realization with decay rate $\beta$.

      %we relate decay of IRR to a finite $\ell_{\infty}$-gain (BIBO stability) and an $\ell_{\infty}$-forgetting property of $H$. 
      %In fact,  
      \textcolor{black}{Next, 
      we show that BIBO stability together with $\ell_{\infty}$-forgetting is equivalent to the existence of a GUES realization.}
  % However, BIBO stability alone is not sufficient for the existence of a GUES realization of $H$.
  %The following example shows that this is not the case. Consider a response map $H$ defined as follows: for any $w=(p(0),v(0)) \cdots (p(t),v(t)) \in (\mathbb{P} \times \mathbb{R})^{+}$, where $\mathbb{P}=\{e_1,e_2\}$, $H(w)=\sum_{s=0}^{t-1} p_1(s) v(s)$. It can be easily verified that $H$ is BIBO stable, but it does not have a GUES realization. In fact, the impulse response of $H$ for the scheduling sequence $p$ such that $p(s)=e_1$ for all $s$ is constant and does not decay to zero.
  We say that $H$ has the \emph{$\ell_{\infty}$-forgetting property} if there exists a decreasing sequence $\{G_{H,k}\}_{k=0}^{\infty}$ 
  \textcolor{black}{of real numbers}
  \textcolor{black}{such that 
  \textbf{(1)} $\lim_{k \rightarrow \infty} G_{H,k}=0$ and
  %\textcolor{black}{with the property that} 
  \textbf{(2)} for any
  scheduling signal $\underline{p} \in (\mathbb{P})^{t}$, $t \ge 1$ and 
  any
  two input signals $\underline{v}_1,\underline{v}_2  \in (\mathbb{R}^{m})^{t}$
 %  \in (\mathbb{P} \times \mathbb{R}^{m})^{+}$,
  $i=1,2$ % such that the scheduling signal is the same, and the input signals 
  such that $\underline{v}_1$ and $\underline{v}_2$} are the same for the last \textcolor{black}{$0 \le k \le t$} time steps, i.e., $\underline{v}_1(s)=\underline{v}_2(s)$ for $s=t-k+1,\ldots t$, \textcolor{black}{ the following holds with 
  $w_i=((\underline{p}[1],\underline{v}_i[1]) \cdots (\underline{p}[t],\underline{v}_i[t]))$, 
  $i=1,2$:}
  \begin{align}
  & \|H(w_1)\|_2 \le G_{H,0} \|\textcolor{black}{\underline{v}_1}\|_{\infty} \label{eq:l_infty_gain1} \textcolor{black}{,} \\
  & \|H(w_1) - H(w_2)\|_2 \le G_{H,k} \|\textcolor{black}{\underline{v}_1}-\textcolor{black}{\underline{v}_2}\|_{\infty}
  \label{eq:l_infty_gain2} \textcolor{black}{.}
  \end{align}
  We say that $H$ is \emph{exponentially  $\ell_{\infty}$-forgetting} \textcolor{black}{with rate $\beta \in (0,1)$}, if, \textcolor{black}{in addition to \eqref{eq:l_infty_gain1}-\eqref{eq:l_infty_gain2},} $G_{H,k}<\textcolor{black}{K}\beta^k$ \textcolor{black}{holds} for some $\textcolor{black}{K} > 0$ and all $k \in \mathbb{N}$. 
    \textcolor{black}{%This definition encodes uniform BIBO stability and input forgetting for all scheduling sequences:
    Intuitively, \eqref{eq:l_infty_gain1} ensures BIBO stability by giving a uniform $\ell_\infty$-gain bound, while \eqref{eq:l_infty_gain2} ensures that the effect of inputs more than $k$ steps in the past decays uniformly. Note that \eqref{eq:l_infty_gain1} alone implies absolute summability and forgetting for each fixed scheduling, but not uniformly across all scheduling signals \cite[Chapter~7d, Thm.~17]{Cal:Des}.}
  %this already implies an input forgetting property, i.e., 
  %a version of \eqref{eq:l_infty_gain2} where $G_{H,k}$ is replaced by a constant which depends on the scheduling signal. 
  %We need \eqref{eq:l_infty_gain2} to ensure that the rate with which past inputs are forgotten
  %is independent of the scheduling. 
  %Clearly, $G_{H,0}$ is the $\ell_{\infty}$-gain of $H$, and hence the $\ell_{\infty}$-forgetting %property implies BIBO stability of $H$.
  %By viewing LPV systems as time-varying linear systems,
  %By \cite[Chapter 7d, Theorem 17]{Cal:Des}, 
  %\eqref{eq:l_infty_gain1} implies absolute summability of the impulse response along each scheduling signal, which implies a weaker  version of \eqref{eq:l_infty_gain2}


  %Next, we state the main theorem that relates the $\ell_{\infty}$-forgetting property to the uniform decay of the response maps $\Psi_H$.
  \begin{Theorem}
  If $H$ is IRR and has a finite Hankel-rank, then the following are equivalent: \textbf{(1)} $H$ has the $\ell_{\infty}$-forgetting property, \textbf{(2)} $H$ is exponentially $\ell_{\infty}$-forgetting, \textbf{(3)} $H$ is IRR-UED, \textbf{(4)} $H$ has a realization by a GUES LPV-SSA/LSS. 
  \textcolor{black}{Moreover, }if $H$ \textcolor{black}{is} $\ell_{\infty}$-forgetting, then 
  it \textcolor{black}{is} exponentially $\ell_{\infty}$-forgetting with a rate $\beta$,
  and % if and only if all minimal LPV-SSA realizations of $H$ are GUES with decay rate $\beta$.
  all minimal LPV-SSA/LSS realizations of $H$ are GUES with \textcolor{black}{the rate $\beta$}. 
  %\begin{enumerate}
  %\item $H$ has the $\ell_{\infty}$-forgetting property $\implies$ $H$ is IRR-UD.
  %\item $H$ is IRR-UAS $\implies$ $H$ has the $\ell_{\infty}$-forgetting property.
  %\item $H$ is IRR-UD $\iff$  $h \diamond \underline{p}$ is uniformly decaying for all $\underline{p} \in \mathbb{P}^{\infty}$, i.e., $\lim_{N \rightarrow +\infty} \|(h \diamond \underline{p})(N+\tau,\tau)\|_2=0$.
  %\item $H$ is IRR-UED $\iff$  the impulse response $h \diamond \underline{p}$ is uniformly exponentially decaying for all $\underline{p} \in \mathbb{P}^{\infty}$, i.e., there exists $\beta \in (0,1)$ such that $\sup_{\underline{p} \in \mathbb{P}^{\infty}} \|(h \diamond \underline{p})(N+\tau,\tau)\|_2 < C\beta^N$ for some $C > 0$.
  %\item $H$ is IRR-UAS $\iff$  the impulse response $h \diamond \underline{p}$ is uniformly absolutely summable for all $\underline{p} \in \mathbb{P}^{\infty}$, i.e., $\sup_{\underline{p} \in \mathbb{P}^{\infty}} \sum_{r=0}^{N-1} \|(h \diamond \underline{p})(N,r)\|_2 < +\infty$.
  %\end{enumerate}
  \end{Theorem}
  \begin{proof}
    The \textcolor{black}{theorem follows} from Corollary \ref{cor:lpv} and the following implications:

    \textbf{$\ell_{\infty}$-forgetting property $\implies$ $H$ is IRR-UD.}
      For \textcolor{black}{a} $\underline{p} \in \mathbb{P}^{\infty}$, \textcolor{black}{$k \ge 1$},
      define the response map
        \( H_{\underline{p},k}:(\mathbb{R}^m)^{+} \to \mathbb{R}^{\textcolor{black}{n_y}}\)
        as follows: for any $\underline{v} \in (\mathbb{R}^m)^{t}$,
        $t \ge 1$,
          let $w \in (\mathbb{P} \times \mathbb{R}^m)^{t+k}$ be such that
          $w[\tau]=(\underline{p}[\tau],\underline{v}[\tau])$ for $\tau=1,\ldots, t$, and $w[\tau]=(\underline{p}[\tau],0)$ for $\tau=t+1,\ldots, t+k$. Then we set
          %((\underline{p}[1],\underline{v}[1]) \cdots (\underline{p}[t],\underline{v}[t])(\underline{p}[t+1],0)\cdots (\underline{p}[t+k],0))$, and set
          %\begin{equation}  
           % \label{eq:H_p_k}
          $H_{\underline{p},k}(\underline{v}[1] \cdots \underline{v}[t]) = H(w)$. 
          %\textcolor{black}{and set $H_{\underline{p},0}(\underline{v}[1] \cdots \underline{v}[t]) = H((\underline{p}[1],\underline{v}[1]) \cdots (\underline{p}[t],\underline{v}[t]))$.}
          %\end{equation}
          It follows that 
          \textcolor{black}{for any sequences $\underline{v}_1,\underline{v}_2 \in (\mathbb{R}^m)^{t+k}$ such that 
          $\underline{v}_1[s]-\underline{v}_2[s]=\left\{\begin{array}{rl} \underline{v}[s] & s \le t \\ 0 & s > t \end{array}\right.$}
          \begin{equation}
            \label{eq:F_p_k}
          %\begin{split}
             H_{\underline{p},k}(\underline{v})
            \!\!=\!\! \sum_{\tau=1}^{t} (h \diamond \underline{p})(t+k,\tau)\underline{v}[\tau]=H(w_1) \!\!-\!\! H(w_2) \textcolor{black}{,}
          %\end{split}  
            \end{equation}
          \textcolor{black}{where $w_i=(\underline{p}[1],\underline{v}_i[1]) \cdots (\underline{p}[t+k],\underline{v}_i[t+k])$,}
          $i=1,2$.
          %for $\tau=1,\ldots, t$, and $w_2[\tau]=(\underline{p}[\tau],0)$ for $\tau=1,\ldots, t$, and $w_i[\tau]=(\underline{p}[\tau],0)$ for $\tau=t+1,\ldots, t+k$
          %. That is, $w_1
          %=((\underline{p}[1],\underline{v}[1]) \cdots (\underline{p}[t+k],\underline{v}[t+k])))$ and 
          %$w_2=((\underline{p}[1],0) \cdots (\underline{p}[t],0)(\underline{p}[t+1],\underline{v}[t+1]) \cdots (\underline{p}[t+k],\underline{v}[t+k])))$.
          %When $k=0$, we set 

          \begin{color}{black}
          For any output index $r=1,\ldots,n_y$ let $\underline{v}^{\textcolor{black}{r}} \in (\mathbb{R}^m)^t$ be such that 
          $\textcolor{black}{(\underline{v}^r[\tau])_j}=\left\{\begin{array}{rl} 1 & ((h \diamond \underline{p})(t+k,\tau))_{r,j}  > 0\\ -1 & \mbox{otherwise} \end{array}\right. $ is the sign of the $(r,j)$th entry of $(h \diamond \underline{p})(t+k,\tau)$ for
          $j=1,\ldots,m$.
          % (the sign of a zero is understood to be zero). 
          Then, 
         %$(h \diamond \underline{p})(t+k,\tau)$, we can make the $r$th entry of
         \( (H_{\underline{p},k}(\underline{v}^{\textcolor{black}{r}}))_r=\sum_{\tau=1}^{t} \sum_{j=1}^{m} |\left((h \diamond \underline{p})(t+k,\tau)\right)_{r,j}| \).  % \] \le G_{H,k}. \]
         %Therefore, 
         %$\sum_{\tau=1}^{t}\|(h \diamond \underline{p})(t+k,\tau)\|_2 \le  \le  pG_{H,k}$
         %for some sequence $\{G_{H,k}\}$ with $\lim_{k \rightarrow +\infty} G_{H,k}=0$.
         % property is equivalent to
         %existence of a decreasing sequence $\{G_{H,k}\}$ such that 
         %$\lim_{k \rightarrow +\infty} G_{H,k}=0$  such that
         %\begin{equation}
          %  \label{eq:F_p_k:2}
           % \begin{split}
           Hence 
          as $H$ is  $\ell_{\infty}$-forgetting, 
          \textcolor{black}{by setting $\underline{v}_2=0$ and $\underline{v}_1[s]=\underline{v}^{\textcolor{black}{r}}[s]$ for $s=1,\ldots,t$, it follows that} 
          $\textcolor{black}{(H_{\underline{p},k}(\underline{v}^{\textcolor{black}{r}}))_r \le} \|H_{\underline{p},k}(\underline{v}^{\textcolor{black}{r}})\|_{2}=\|H(w_1)-H(w_2)\|_2 \le G_{H,k}\|\underline{v}^{\textcolor{black}{r}}\|_{\infty} \textcolor{black}{\le G_{H,k}}$, and thus
          \begin{align*}
           & \sum_{\tau=1}^{t} \|(h \diamond \underline{p})(t+k,\tau)\|_2 \le  
           \sum_{\tau=1}^{t} \|(h \diamond \underline{p})(t+k,\tau)\|_F 
           \\
           & \le \sum_{\tau=1}^{t} \sum_{r=1}^{\textcolor{black}{n_y}} \sum_{j=1}^{m} |\left((h \diamond \underline{p})(t+k,\tau)\right)_{r,j}| \le 
         n_y G_{H,k} \textcolor{black}{.}
          \end{align*}
          %\end{split}
         %\end{equation}
         As $\lim_{k \rightarrow \infty} G_{H,k}=0$,
         it then follows that $(h \diamond \underline{p})$ is IRR-UD. 
        \end{color}
         %Since the Euclidean norm of $H_{\underline{p},k}(\underline{v})$ is not smaller than 
         %its $r$-th component, it follows that 
         %$\sum_{\tau=1}^{t} \|(h \diamond \underline{p})_r(t+k,\tau)\|_1 \le pG_{H,k}$ 
         %for any $r=1,\ldots p$, 
         %and hence
         %and $\tau$. 
         %In particular, 
         %and $t \ge \tau$, and 
         %uniformly decaying for all $\underline{p} \in \mathbb{P}^{\infty}$, and hence $H$ is IRR-UD.
        % it follows that for any $\epsilon > 0$, there exists $K_{\epsilon}$ such that for any $k > K_{\epsilon}$, and any $N$ and $\tau$, $\|(h \diamond \underline{p})(N+k,\tau)\|_2 < \epsilon$,
         %and hence $\lim_{k \rightarrow +\infty} \sup_{N} \|(h \diamond \underline{p})(N+k,\tau)\|_2=0$.
         %$G_{H,k}$ is the $\ell_{\infty}$-gain of $H_{\underline{p},k}$.
         % However, as $h_k \diamond p$ is the impulse response of $H_{\underline{p},k}$,
        %from \cite[Chapter 7d, Theorem 17]{Cal:Des} it follows that
        %if $H_{\underline{p},k}$ has finite $\ell_{\infty}$-gain $G_{H,k}$, then 
         % \( \sup_{N} \sum_{r=0}^{N} \|(h \diamond \underline{p})(N+k,r)\|_2 \le G_{H,k} \).
          % is finite, and 
          %in fact the latter is smallest possible $\ell_{\infty}$-gain of $H_{\underline{p},k}$
         % as for a suitable choice of $\underline{v} \in ([-1,1]^m)^{N}$, 
         % $\sum_{r=0}^{N-1} \|(h \diamond \underline{p})(N+k,r)\|_2=H_{\underline{p},k}(\underline{v})$. 
         %By virtue of \eqref{eq:pf:lpv1}, the $\ell_{\infty}$-forgetting property is equivalent to the existence of a decreasing sequence $\{G_{H,k}\}$ such that $\lim_{k \rightarrow +\infty} G_{H,k}=0$ and
        %\[ 
        %  \sup_{q_0,\ldots,,q_N \in Q, N \ge 1} \sum_{r=0}^{N-1} \|S_{q_{N+k}, q_r}(q_{r+1} \cdots q_{N+k-1})\|_2 \le G_{H,k}
        %\] 
    \textbf{IRR-UED with rate $\beta$ $\implies$ $H$ exponential $\ell_{\infty}$-forgetting with rate $\beta$.}
         If $H$ is IRR-UED, then set % for all $k \ge 0$,
         \[
		 G_{H,k}=\textcolor{red}{\sqrt{m}}\sup_{t \ge 1, \underline{p} \in \mathbb{P}^{\infty}} \sum_{\tau=1}^{\textcolor{black}{\min\{t,t+k-1\}}} \|(h \diamond \underline{p})(t+k,\tau)\|_2
         \]
         for all $k \ge 0$.
         \textcolor{black}{Since $H$ is IRR-UED}, it follows that 
         $\|(h \diamond \underline{p})(t+k,\tau)\|_2 \le C \beta^{t+k-\tau\textcolor{black}{-1}}$ for some $C > 0$ and $\beta \in (0,1)$, and hence $G_{H,k} \le K\beta^k$ for 
	 $K=\textcolor{red}{\sqrt{m}}C/\textcolor{black}{(\beta(1-\beta))}$.
         Recall the definition of $H_{\underline{p},k}$ from the previous implication.
         \textcolor{black}{By repeating the argument of \cite[Chapter 7d, Theorem 17]{Cal:Des},}
         it follows that 
         \textcolor{black}{$G_{H,0}$ satisfies \eqref{eq:l_infty_gain1}, and}
         the $\ell_{\infty}$-gain of $H_{\underline{p},k}$ is upper bounded by $G_{H,k}$, \textcolor{black}{$k \ge 1$,} \textcolor{black}{i.e., $\|H_{\underline{p},k}(\underline{v})\|_2 \le G_{H,k} \|\underline{v}\|_{\infty}$.}
         Then \eqref{eq:F_p_k} implies that $G_{H,k}$ satisfies \eqref{eq:l_infty_gain2},
         for $k \ge 1$. %and from the definition of $H_{\underline{p},0}$ 
         %it follows that 
         %That is, $H$ is exponentially $\ell_{\infty}$-forgetting with a rate $\beta$.
         %and hence $H$ has the $\ell_{\infty}$-forgetting property.
         %In particular, for $k=0$, $G_{H,0}$ is an upper bound on the $\ell_{\infty}$-gain of $H$.
         %By virtue of $H$ being IRR-UAS, $G_{H,k}$ is finite for all $k$ and $\lim_{k \rightarrow +\infty} G_{H,k}=0$. Finally, from \eqref{eq:F_p_k}
         %µit follows that $G_{H,k}$ satisfies \eqref{eq:l_infty_gain2}, and hence $H$ has the $\ell_{\infty}$-forgetting property.
         %By virtue of \eqref{eq:F_p_k}
       %  -\eqref{eq:H_p_k} it follows that for any $\underline{p} \in \mathbb{P}^{\infty}$, $H_{\underline{p},k}$ has $\ell_{\infty}$-gain upper bounded by $G_{H,k}$, and hence the $\ell_{\infty}$-gain of $H_{\underline{p},k}$ is upper bounded by $G_{H,k}$.
        % equal to $\sup_{N \ge 0} \sum_{r=0}^{N-1} \|(h \diamond \underline{p})(N+k,r)\|_2$ for any $\underline{p} \in \mathbb{P}^{\infty}$, and hence by taking the supremum over all $\underline{p} \in \mathbb{P}^{\infty}$, we obtain that
  %
  %
  %           \sum_{r=0}^{N-1} \|h \diamond \underline{p}(N+k,r)\|_2 = \sup_{N \ge 0, q_0,\ldots, q_{N+k-1} \in Q}
   %      \]     
   %      it follows that 
   %      \[
   %        \begin{split}
   %         & \sum_{r=0}^{N-1} \|S_{q_{N+k},q_r}(q_{r+1}\cdots q_{N+k-1})\|_2 \le   \\
   %         m \sup_{\substack{q \in Q\\ j=1,\ldots,m}} 
   %          \sum_{r=0}^{N-1} \|M_{q,}(q_{r+1}\cdots q_{N+k-1})\|_2 < K=:G_{H,0}
   %        \end{split}
   %      \]
   %      Moreover, for every $k > 0$, let 
   %      \[
   %      \begin{split}
   %       & G_{H,k}=m \!\!\! \sup_{\substack{q_1,\ldots, q_{N+k-1} \in Q\\ (q,j) \in Q \times \{1,\ldots,m\},\, N > 0}}
   %            \sum_{r=0}^{N-1} \|M_{q,j}(q_{r+1}\cdots q_{k+N-1})\|_2
   %      \end{split}
  %       \]
  \end{proof}
  %From Theorem \ref{theo:main} applied to $\Psi_H$ and the correspondence between LPV-SSAs and SLSs, we obtain the following results.
  %\begin{corollary}
  %If $\Psi_H$ satisfies the finite Hankel-rank condition, then the following properties of $H$ are equivalent: IRR-UD, IRR-UAS, IRR-UED, and the $\ell_{\infty}$-forgetting property.
  %µA response map $H$ can be realized by a GUES LPV system if and only if it has an impulse response representation (IRR), it satisfies the finite Hankel-rank condition, and it is either IRR-UD, IRR-UED, IRR-UAS, or has the $\ell_{\infty}$-forgetting property.
  %\end{corollary}
  % That is, if $H$ can be realized by a LPV-SSA/LSS, then relatively mild decay properties (UD) of the impulse response of $H$ are equivalent to the existence of a GUES realization of $H$ and to exponential decay and uniform summability properties. Moreover, the latter is equivalent to forgetting past inputs at a rate which is proportional to its amplitude ($\ell_{\infty}$-forgetting property).


\section{Conclusions}
We characterized when an input--output family admits a stable finite-dimensional realization (SLS, SAS, LPV-SSA, LSS): for finite Hankel-rank, stability is equivalent to uniform decay/forgetting, and minimal realizations inherit the same decay rate.
A future direction is to develop the relationship with echo-state/fading-memory properties and input--output tests for existence of quadratic/polyhedral/SOS Lyapunov functions.

\textbf{Acknowledgements:} we thank the anonymous reviewers for numerous useful technical remarks and suggestions. 

 \bibliographystyle{IEEEtran}
\bibliography{petreczky_chapterV2}
\end{document}
